\documentclass[11pt]{article}
\usepackage[margin=0.92in]{geometry}
\usepackage[T1]{fontenc}
\usepackage{lmodern}
\usepackage{microtype}
\usepackage{amsmath,amssymb,amsthm,mathtools}
\usepackage{booktabs,longtable,array,tabularx}
\usepackage{enumitem}
\usepackage{xcolor}
\usepackage{listings}
\usepackage{xurl}
\usepackage{hyperref}
\usepackage[nameinlink,capitalise,noabbrev]{cleveref}
\usepackage{fancyhdr}
\usepackage{caption}

\definecolor{deepblue}{RGB}{22,55,92}
\definecolor{softgray}{RGB}{245,246,248}
\definecolor{codegreen}{RGB}{18,105,58}
\hypersetup{
  colorlinks=true,
  linkcolor=deepblue,
  citecolor=deepblue,
  urlcolor=deepblue,
  pdftitle={OCiPRA: Observation-Contract Internalization for iPRA},
  pdfauthor={Prepared for Vladimir Reshetnikov}
}
\pagestyle{fancy}
\fancyhf{}
\lhead{OCiPRA}
\rhead{15 November 2026}
\cfoot{\thepage}
\setlength{\headheight}{14pt}
\setlength{\parindent}{1.15em}
\setlength{\parskip}{0.25em}
\setlist{itemsep=0.2em,topsep=0.35em}
\emergencystretch=2em

\lstdefinelanguage{Lean}{
  morekeywords={def,theorem,structure,inductive,namespace,end,where,by,exact,apply,have,let,import,open,variable,forall,exists,match,with},
  sensitive=true,
  morecomment=[l]{--},
  morecomment=[s]{/-}{-/},
  morestring=[b]"
}
\lstset{
  basicstyle=\ttfamily\small,
  keywordstyle=\bfseries\color{deepblue},
  commentstyle=\itshape\color{codegreen},
  backgroundcolor=\color{softgray},
  frame=single,
  framerule=0.2pt,
  breaklines=true,
  columns=fullflexible,
  keepspaces=true,
  showstringspaces=false
}

\newtheorem{theorem}{Theorem}[section]
\newtheorem{lemma}[theorem]{Lemma}
\newtheorem{proposition}[theorem]{Proposition}
\newtheorem{corollary}[theorem]{Corollary}
\newtheorem{principle}[theorem]{Principle}
\theoremstyle{definition}
\newtheorem{definition}[theorem]{Definition}
\newtheorem{example}[theorem]{Example}
\newtheorem{warning}[theorem]{Warning}
\newtheorem{obligation}[theorem]{Formalization obligation}
\theoremstyle{remark}
\newtheorem{remark}[theorem]{Remark}

\newcommand{\iPRA}{\mathsf{iPRA}}
\newcommand{\PRAAx}{\mathsf{PRAAx}}
\newcommand{\NPrf}{\mathsf{NPrf}}
\newcommand{\Prf}{\mathsf{Prf}}
\newcommand{\Nat}{\mathbb N}
\newcommand{\Bool}{\mathbb B}
\newcommand{\code}[1]{\texttt{#1}}
\newcommand{\Obs}{\mathsf{Obs}}
\newcommand{\Valid}{\mathsf{Valid}}
\newcommand{\CertCheck}{\mathsf{Check}}
\newcommand{\Good}{\mathsf{Good}}
\newcommand{\diam}{\operatorname{diam}}
\newcommand{\pair}[2]{\langle #1,#2\rangle}
\newcommand{\repo}{\textsc{ProveIt}}

\title{\textbf{OCiPRA: Observation--Contract Internalization\\for Intuitionistic Primitive Recursive Arithmetic}\\[0.5em]
\large A proof-producing framework for finite, relational, approximate, and conditional mathematics}
\author{Research report prepared for Vladimir Reshetnikov}
\date{15 November 2026 (America/Los\_Angeles)}

\begin{document}
\maketitle

\begin{abstract}
The current \repo{} iPRA collection contains a broad family of theorem packets and several increasingly capable finitary frameworks.  The recurring difficulty is no longer merely coding mathematics by natural numbers.  It is keeping four boundaries explicit at the same time: finite observation versus an infinite ideal object, verifier versus producer, semantic correctness versus an object-language derivation, and deterministic output versus a genuinely relational choice.  This report develops \emph{Observation--Contract Internalization for iPRA} (OCiPRA), a common framework for those boundaries.

An OCiPRA construction is a typed finite contract.  It names the primitive-recursive programs used at each node, the reverse observation demand, a primitive-recursive resource bound, a bounded verifier, and local proofs whose conclusions are actual \(\NPrf\) or \(\Prf\,\PRAAx\) judgments.  A fixed contract DAG is compiled by a host-level proof-producing recursion.  Equation-heavy nodes use the repository's native scalar calculus and are transported by \code{native\_iff\_pra}; bounded logic and witness packaging use the full \(\Prf\,\PRAAx\) calculus.  No universal unclocked evaluator and no semantic reflection principle is assumed.

The report proves, in full human-readable detail, bounded-formula Booleanization, least bounded search, deterministic contract-DAG composition, clocked trace reflection, monotone finite demand closure, finite relational viability, exact branch elimination under observable invariance, approximate elimination under a diameter certificate, and the producer--verifier--derivation factorization.  It then applies these results to the repository's reported obstacles in analysis, algebra, geometry, topology, probability, Diophantine geometry, descent, and prime distribution.  The accompanying exact-arithmetic Python reference implementation has eleven passing tests.  The Lean files give definitions, actual object-language interfaces, a small existential example, exact target signatures, and a dependency-ordered completion path.  Lean was not available in the artifact runtime, so no new kernel-checked application theorem is claimed.
\end{abstract}

\tableofcontents
\newpage

\section{Status, scope, and repository baseline}

\subsection{What is being extended}

At the pinned repository commit, the primitive-recursive arithmetic directory distinguishes two bodies of work.  The \code{ConciseIPRA} library is a Lean development of mutually bi-interpretable presentations of intuitionistic primitive recursive arithmetic.  The proof-packet tree contains seventy-two delivered research packets.  The repository's own status page says that six packets construct some \(\Prf\,\PRAAx\) derivations for framework rules or small examples, but no packet claims such a derivation for its main application theorem.  This distinction is the starting point of the present report; it is not a criticism of the substantial human mathematics in the packets.

Several foundation packets already provide essential pieces:

\begin{itemize}
  \item finite observations, finite contracts, and clocked trace reflection;
  \item backward precision demand and changing-parameter recursion;
  \item finite relational clouds, viability, and stable observables;
  \item complete bounded enumeration from supplied height or algebra certificates;
  \item small, actual object-language derivations for composition, induction, iteration, simulation, and observation rules;
  \item a current bridge between the native scalar calculus and ordinary iPRA.
\end{itemize}

OCiPRA does not rename these pieces and present them as new.  It adds a single statement calculus in which they fit, identifies the exact object-proof inputs required by a global compiler, supplies a dual proof backend, proves a finite closure theorem for cyclic demand, and turns the recurring caveats into a mechanical admissibility test.

\subsection{The object theory used here}

The intended object theory is the repository theory \(\PRAAx\): first-order intuitionistic logic with equality, a function symbol for every fixed primitive-recursive definition, defining equations for projection, composition, and primitive recursion, successor axioms, and induction for quantifier-free formulas.  We write
\[
   \Gamma\vdash_{\PRAAx}\varphi
\]
for the repository judgment represented in Lean by \code{Prf PRAAx}.  This notation is deliberately different from ``Lean proves a proposition about natural numbers.''

The concise native calculus has iterator terms and formulas generated by equations and the two binders \(Q\) and \(D\).  Its derivability judgment is \(\NPrf\).  The existing theorem
\[
  \NPrf\;\Gamma\;A
  \quad\Longleftrightarrow\quad
  \Gamma^{\mathrm{PRA}}\vdash_{\PRAAx} A^{\mathrm{PRA}}
\]
is implemented as \code{native\_iff\_pra}.  This theorem gives a strong and economical backend for generated equational proofs.  It does \emph{not} say that a true semantic assertion automatically has a proof.

\subsection{Four distinctions that must never be collapsed}

\begin{definition}[Observation versus ideal object]
A finite observation is a natural-number code that can be produced and checked at a specified precision or query.  An ideal object---a real number, infinite series, branch, compact space, measure, or quotient object---is represented only through such observations and coherence laws.  It is not silently added as a higher-type object of iPRA.
\end{definition}

\begin{definition}[Producer versus verifier]
A \emph{producer} supplies a certificate, bound, finite cover, modulus, representative, or exceptional datum.  A \emph{verifier} checks it by a primitive-recursive bounded computation.  The conditional verifier theorem may be internal to iPRA even when no uniform producer theorem is known.
\end{definition}

\begin{definition}[Semantic theorem versus object derivation]
A host theorem establishes a mathematical property in the metatheory.  An object derivation is a value of \(\Prf\,\PRAAx\) or \(\NPrf\).  A compiler may use a host proof to justify its implementation, but its output must still have the object-derivation type.
\end{definition}

\begin{definition}[Selector versus relation]
A selector returns one exact representative.  A relation returns a finite family of acceptable representatives.  A selector is justified only by a bounded search theorem or by a stability/isolation hypothesis that makes the selected observable independent of the branch.
\end{definition}

Many failed or incomplete formalizations can be diagnosed as one of these four collapses.  For example, a proof of finiteness is treated as if it produced a stopping bound; a passing numerical checker is treated as if it were a \(\Prf\) term; a root cluster is treated as if it supplied a continuously labelled root; or an external analytic constant is omitted from the theorem statement.

\subsection{Exact contribution and exact non-claims}

The main mathematical contribution is the compilation theorem of \cref{thm:master-compiler}.  It is a metatheorem about constructing object proofs from local object proofs.  Its proof is constructive and finite.  The supplied Python implementation exercises the computational content but is not the proof object promised by that theorem.  The supplied Lean outline isolates the remaining engineering steps.

This report does not claim:

\begin{itemize}
  \item a newly kernel-checked generic compiler;
  \item a general effective form of Faltings' theorem, Siegel's theorem, uniform Mordell--Weil generation, Chebotarev, Siegel--Walfisz, or Bombieri--Vinogradov;
  \item an exact Puiseux or eigenvector selector through collisions;
  \item an unclocked universal evaluator for arbitrary primitive-recursive program codes;
  \item a verification of the seventy-two proof packets;
  \item that every classical theorem has an iPRA reformulation with the same surface statement.
\end{itemize}

The appropriate standard is stronger than ``the idea is finitary'' and weaker than ``all mathematics is already formalized.''  Every theorem below states its finite data, its resource bound, the status of any producer, and the intended object-proof endpoint.

\section{The OCiPRA statement calculus}

\subsection{Finite observation presentations}

\begin{definition}[Deterministic observation presentation]
A deterministic observation presentation \(X\) consists of fixed primitive-recursive data
\[
  \mathsf{Name}_X,\quad \mathsf{Query}_X,\quad
  \Obs_X(n,q),\quad \Valid_X(n,q,a),
\]
together with a primitive-recursive refinement relation \(q'\succeq q\), restriction maps, and object-language proofs of the following local laws:
\begin{enumerate}[label=(\roman*)]
  \item \(\Valid_X(n,q,\Obs_X(n,q))\);
  \item restriction of a valid finer observation is a valid coarser observation;
  \item two valid observations of the same name and query agree in the chosen observable sense.
\end{enumerate}
All values are natural-number codes.  ``Name'' is therefore a role played by a number and fixed programs, not a new sort in the object language.
\end{definition}

\begin{definition}[Relational observation presentation]
A relational presentation replaces the single output by a bounded relation
\[
  R_X(n,q,a),\qquad a\leq B_X(n,q),
\]
and proves bounded totality.  Coherence is expressed by a bounded compatibility relation between answers at related queries.  It is not part of the definition that one answer be canonically distinguished.
\end{definition}

\begin{definition}[Stable observable]
Let \(R_X(n,q,a)\) be a relational presentation.  An exact stable observable is a primitive-recursive map \(O(a)\) with an object proof
\[
 R_X(n,q,a)\wedge R_X(n,q,b)\;\longrightarrow\;O(a)=O(b).
\]
A quantitative stable observable has a rational-valued code \(O(a)\), a requested precision \(p\), and a proof that every two valid outputs have distance at most \(2^{-p}\), or another explicitly coded error budget.
\end{definition}

This formulation contains exact finite objects as a degenerate case: there is one query and one exact observation.  It contains regular Cauchy names, jets, finite nets, rational distributions, polynomial truncations, root clusters, spectral lists, and finite covers by changing only the query type and coherence law.

\subsection{Proof-carrying local contracts}

A deterministic node with \(n\) numerical inputs is represented by:

\begin{enumerate}
  \item a precondition \(P(\vec x)\);
  \item a fixed primitive-recursive program \(f(\vec x)\);
  \item a postcondition \(Q(y,\vec x)\);
  \item an object derivation
  \[
       P(\vec x)\vdash_{\PRAAx} Q(f(\vec x),\vec x).
  \]
\end{enumerate}

This is exactly the logical content of the \code{ProgramContract} structure in the supplied Lean outline.  The existential result follows by the ordinary existential-introduction constructor:
\[
  P(\vec x)\vdash_{\PRAAx}\exists y\,Q(y,\vec x).
\]
No soundness reflection is used; the proof field is already a derivation.

For a relational node the contract contains directly
\[
  P(\vec x)\vdash_{\PRAAx}\exists y\,R(y,\vec x).
\]
A later bounded-search compiler may turn this relation into a canonical least witness, but only when the relation and bound are explicit.  Otherwise the relational theorem is the honest endpoint.

\subsection{Resource and demand effects}

A contract carries two numerical effects.

\begin{definition}[Demand]
The demand transformer \(\delta_f(p)\) states how precise each input observation must be in order to produce output precision \(p\).  For several inputs it is a finite vector.  A node with several downstream consumers receives the componentwise maximum of their demands.
\end{definition}

\begin{definition}[Resource]
The resource bound \(\rho_f(\vec x,p)\) is a primitive-recursive bound on all loops, trace lengths, table sizes, candidate ranges, or descent ranks used by the verifier and proof.  It is part of the theorem data, not an afterthought about implementation efficiency.
\end{definition}

The corresponding informal typing judgment is
\[
   \Gamma\vdash
   f : X_1\times\cdots\times X_k
   \xRightarrow[\rho_f]{\delta_f}
   Y
   \;\bigm|\; Q_f .
\]
It means that the fixed program \(f\), when supplied observations meeting \(\delta_f\), runs within \(\rho_f\) and satisfies the postcondition \(Q_f\), with an actual local object derivation.

\subsection{The four theorem classes}

\begin{center}
\begin{tabularx}{\textwidth}{@{}p{0.19\textwidth}X X@{}}
\toprule
Class & Normal form & Typical output \\
\midrule
Deterministic & fixed PR program plus local \(\Prf\) proof & exact finite calculation, regular-name operation \\
Bounded relational & bounded total relation plus viability/search proof & root cluster, finite cover, candidate witness \\
Stable observable & relational output plus exact invariance or diameter proof & ordered eigenvalue multiset, quotient value, robust chart datum \\
Conditional producer & \(\CertCheck(x,c)\to\mathsf{Goal}(x,c)\), with existence of \(c\) separate & height-bound verifier, analytic-constant verifier, descent representative \\
\bottomrule
\end{tabularx}
\end{center}

Every proposed theorem should first be assigned one of these classes.  A large portion of formalization difficulty comes from insisting that a bounded relational or conditional-producer theorem be stated as an unconditional deterministic theorem.

\subsection{Admissibility test}

A theorem is ready for object internalization when all of the following are explicit:

\begin{enumerate}
  \item every program is fixed at compile time, or an input-coded computation has a fixed primitive-recursive clock and trace verifier;
  \item every loop and search has a primitive-recursive bound;
  \item every certificate predicate is bounded and decidable;
  \item every local semantic correctness statement has been converted to \(\NPrf\) or \(\Prf\,\PRAAx\);
  \item every external producer is either proved, supplied as an argument, or visibly left as a conditional hypothesis;
  \item every exact selector is protected by bounded least search, branch invariance, or a separation gap;
  \item every use of compactness, convergence, or finiteness has been replaced by a modulus, finite net, height bound, complete list, or fixed-precision formulation.
\end{enumerate}

The reference implementation turns this list into a diagnostic classifier.  The classifier is not a proof assistant.  Its role is to reject malformed theorem statements before expensive formalization begins.

\section{Finite coding, bounded formulas, and least search}

\subsection{Primitive-recursive coding of finite data}

The framework uses only first-order numerical objects.  Finite tuples, lists, tagged unions, rational numbers, finite maps, matrices, polynomials, syntax trees, and proof traces are therefore coded by naturals.  Any standard primitive-recursive pairing with primitive-recursive projections is sufficient; the repository already provides pairing operations in its primitive-recursive library.

For definiteness, write \(\pair ab\) for a fixed pairing and \(\pi_0,\pi_1\) for its projections.  A list is coded by
\[
  []=0,\qquad a::s=1+\pair a s.
\]
The functions \(\mathsf{length}\), \(\mathsf{nth}\), append, map for a fixed primitive-recursive function, bounded fold, and bounded membership are primitive recursive.  The proof is by primitive recursion on the list code with an explicit length or fuel bound.  Malformed codes are assigned harmless default values; validity is always a separate bounded predicate.

\begin{lemma}[Finite coding closure]\label{lem:coding}
Every construction obtained from fixed primitive-recursive functions by finite products, finite sums, bounded lists, bounded tables, and bounded folds is primitive recursive.  Its graph is expressible by a quantifier-free formula after adding the corresponding fixed function symbols to iPRA.
\end{lemma}

\begin{proof}
Products and sums are compositions with the pairing and tag functions.  A bounded list operation is primitive recursion on a numerical counter no larger than the supplied length bound.  A bounded table is a list with an index validity predicate.  A bounded fold has an initial state and a fixed primitive-recursive transition; it is exactly primitive recursion.  Since iPRA contains a function symbol and defining equations for each fixed primitive-recursive definition, the graph of the resulting function can be written as an equality between terms.  No unbounded decoding or partial evaluator is involved.
\end{proof}

The word \emph{fixed} matters.  There is no claim that one total primitive-recursive function evaluates every primitive-recursive program code on every input.  A compiler may inspect finite syntax in Lean's metatheory and construct a fixed object-language term.  For input-coded computations it instead verifies a bounded trace with a fixed step function.

\subsection{Bounded formula syntax}

Let \(\mathsf{BFm}_n\) be the formulas over \(n\) free numerical variables generated by:

\begin{itemize}
  \item equalities and inequalities between fixed primitive-recursive terms;
  \item \(\wedge,\vee,\to,\bot\);
  \item \(\exists y<t(\vec x)\,A(y,\vec x)\) and
        \(\forall y<t(\vec x)\,A(y,\vec x)\).
\end{itemize}

Inequalities are merely abbreviations for fixed primitive-recursive comparison tests.  Bounded formulas are intuitionistically decidable because their truth is determined by a finite computation.  The formal point is not merely that a host evaluator exists; the compiler must produce a characteristic primitive-recursive term and an iPRA derivation of equivalence.

\begin{theorem}[Bounded Booleanization]\label{thm:booleanization}
For every fixed bounded formula \(A(\vec x)\) one can effectively construct a primitive-recursive Boolean function \(\chi_A(\vec x)\in\{0,1\}\) and an object-language derivation
\[
   \vdash_{\PRAAx}
   A(\vec x)\leftrightarrow \chi_A(\vec x)=1.
\]
The construction is structural in \(A\).  Its proof uses only the defining equations of the generated functions and quantifier-free induction.
\end{theorem}

\begin{proof}
For an atomic equality \(s=t\), let \(\chi_A\) be the fixed equality-test program.  The predecessor, truncated subtraction, and zero-test programs available in primitive recursive arithmetic give a test satisfying
\[
 \chi_A=1\leftrightarrow s=t,
 \qquad
 \chi_A=0\leftrightarrow s\neq t.
\]
The equivalence has an iPRA proof from the defining equations and the successor axioms.  One may equivalently treat equality tests as a reusable local contract.

Suppose characteristic tests and equivalence derivations have been constructed for \(B\) and \(C\).  Define
\[
\begin{aligned}
 \chi_{B\wedge C}&=\chi_B\cdot\chi_C,\\
 \chi_{B\vee C}&=\max(\chi_B,\chi_C),\\
 \chi_{B\to C}&=\max(1-\chi_B,\chi_C),\\
 \chi_{\bot}&=0.
\end{aligned}
\]
All operations are on Boolean values.  The defining equations give the finite truth tables.  The implication case is intuitionistically sound because the induction hypothesis proves the decidability alternative for \(B\): either \(\chi_B=0\), hence \(B\to C\), or \(\chi_B=1\), in which case \(B\to C\) is equivalent to \(C\).  No unrestricted excluded middle is introduced.

For bounded existential quantification, let \(t=t(\vec x)\) and define the fold
\[
 E(0,\vec x)=0,
 \qquad
 E(k+1,\vec x)=\max(E(k,\vec x),\chi_B(k,\vec x)).
\]
A quantifier-free induction on \(k\) proves
\[
 E(k,\vec x)=1
 \leftrightarrow
 \exists y<k\;B(y,\vec x).
\]
The induction invariant is an equality between primitive-recursive tests together with the already Booleanized bounded relation; it can be packaged as a single Boolean equality.  Substitute \(k=t(\vec x)\).

For bounded universal quantification, use
\[
 U(0,\vec x)=1,
 \qquad
 U(k+1,\vec x)=\min(U(k,\vec x),\chi_B(k,\vec x)),
\]
and prove by quantifier-free induction that \(U(k)=1\) exactly when all \(y<k\) satisfy \(B\).  This completes the structural recursion.  Because the formula is fixed, the metatheoretic recursion creates one fixed primitive-recursive program and one finite object proof.  It is not an internal truth predicate for arbitrary formula codes.
\end{proof}

\begin{remark}[Why all-precision induction is avoided]
A statement such as ``two Cauchy names agree at every precision'' contains an unbounded quantifier and is not itself a quantifier-free induction invariant.  OCiPRA fixes the requested precision \(p\), Booleanizes the finite table needed at that precision, performs quantifier-free induction on the finite computation, and only then universally quantifies \(p\) in the surrounding first-order proof.  This order is essential.
\end{remark}

\subsection{Least bounded search}

Let \(R(\vec x,y)\) be bounded and let \(B(\vec x)\) be primitive recursive.  Define a scan state \((f,w)\), where \(f\in\{0,1\}\) records whether a witness has been found.  Initially \((f,w)=(0,0)\).  At stage \(k\):
\[
  (f,w)\longmapsto
  \begin{cases}
    (1,k),&f=0\text{ and }R(\vec x,k),\\
    (f,w),&\text{otherwise.}
  \end{cases}
\]
Run through \(B(\vec x)+1\) stages.

\begin{theorem}[Least bounded witness]\label{thm:least-search}
There is a primitive-recursive total function \(\mu_{\leq B}R(\vec x)\) and an iPRA derivation of
\[
  \exists y\leq B(\vec x)\,R(\vec x,y)
  \;\longrightarrow\;
  \begin{cases}
  R(\vec x,\mu_{\leq B}R(\vec x)),\\
  \mu_{\leq B}R(\vec x)\leq B(\vec x),\\
  \forall z<\mu_{\leq B}R(\vec x)\;\neg R(\vec x,z).
  \end{cases}
\]
If no witness exists, the function returns the harmless default \(0\) together with a false found-bit.
\end{theorem}

\begin{proof}
Booleanize \(R\).  The scan is primitive recursion on \(k\), with the pair state coded as one natural.  Use the following invariant after \(k\) stages:
\[
\begin{aligned}
 f=0 &\longleftrightarrow \forall z<k\;\neg R(\vec x,z),\\
 f=1 &\longrightarrow
   \bigl(w<k\wedge R(\vec x,w)
   \wedge\forall z<w\;\neg R(\vec x,z)\bigr).
\end{aligned}
\]
Every quantifier is bounded by \(k\), so \cref{thm:booleanization} converts the invariant to an equality of Boolean tests.  The base case is immediate.  For the step, split on the old found-bit and the Boolean value of \(R(\vec x,k)\).  There are four finite cases, each discharged by the scan equations.  Quantifier-free induction yields the invariant at \(B+1\).  Under the existence premise, the first equivalence rules out \(f=0\); the second gives validity and minimality.  The proof constructs an actual derivation once the Booleanization and scan equations are supplied as local object proofs.
\end{proof}

\subsection{Finite enumeration}

The least-search theorem extends immediately to complete bounded enumeration.  Scan every candidate \(y\leq B\), append those satisfying \(R\), and maintain the invariant that the current list contains exactly the satisfying candidates below the current counter, without duplication.  List validity, membership, and equality are bounded tests.  Thus:

\begin{corollary}[Complete bounded list]\label{cor:bounded-list}
For every bounded decidable relation \(R(\vec x,y)\) and primitive-recursive bound \(B(\vec x)\), there is a primitive-recursive list program \(L_R(\vec x)\) and an iPRA proof that
\[
  y\in L_R(\vec x)
  \quad\leftrightarrow\quad
  y\leq B(\vec x)\wedge R(\vec x,y).
\]
\end{corollary}

This corollary is the verifier half of many number-theoretic applications.  Its hypothesis is a \emph{bound}; qualitative finiteness alone is insufficient, as shown in \cref{sec:height-obstruction}.

\section{The proof-producing contract compiler}

\subsection{Two proof backends}

OCiPRA uses two existing proof languages rather than forcing every local argument through one syntax.

\paragraph{Native backend.}
Equation-heavy primitive-recursive reasoning is generated in the native scalar calculus.  Iterator zero and successor equations, transport, and the native \(Q\) and \(D\) rules produce a term of \(\NPrf\).  The existing theorem \code{native\_iff\_pra} transports that term to an ordinary iPRA derivation.  The compiler does not first prove a semantic equation in Lean and then appeal to completeness; it constructs the native derivation.

\paragraph{Full iPRA backend.}
Bounded connectives, existential witnesses, finite relations, and contract packaging are generated directly in \(\Prf\,\PRAAx\).  Booleanized invariants are closed with \code{PRA.qfInd}; witnesses use existential introduction; preconditions are discharged by implication introduction.

A mixed node may prove its computational equations natively, transport them to \(\PRAAx\), and combine them with a bounded logical proof.  This division follows the actual strengths of the repository interfaces.

\begin{principle}[Typed proof production]\label{princ:typed}
The compiler's result type is the desired derivability judgment.  Thus a function
\[
  \mathsf{compile}(D,h):\Prf\,\PRAAx\,\varphi_D
\]
is correct because every constructor branch returns a well-typed proof term.  A separate semantic soundness theorem is useful for the foundational library but is not a substitute for this type.
\end{principle}

\subsection{Fixed DAGs and reverse demand}

Let \(D\) be a finite directed acyclic graph with topological order
\(v_0,\ldots,v_{m-1}\).  Each node has:

\begin{itemize}
  \item predecessors \(\mathsf{pred}(i)\subseteq\{0,\ldots,i-1\}\);
  \item a fixed primitive-recursive program \(f_i\);
  \item a local demand vector \(\delta_i(p)\);
  \item a primitive-recursive local resource bound \(\rho_i\);
  \item a precondition \(P_i\), postcondition \(Q_i\), and local object proof.
\end{itemize}

Given a requested output precision \(p\) at the terminal node, define demands in reverse order.  The terminal demand is \(p\).  Whenever node \(i\) is required at demand \(d_i\), each predecessor \(j\) receives the corresponding component of \(\delta_i(d_i)\).  If several consumers require \(j\), take the maximum.  Because the graph and every \(\delta_i\) are fixed and primitive recursive, the reverse pass is primitive recursive.

The forward pass evaluates the nodes in topological order at the computed demands.  Its global state is a finite coded table of demands, outputs, and local certificate bits.  A single invariant says:

\begin{enumerate}
  \item all processed nodes have valid outputs;
  \item the stored demand of each processed dependency is at least the demand requested by every processed consumer;
  \item every stored trace fragment is within its local resource bound;
  \item unprocessed table entries retain their initialization values.
\end{enumerate}

For a fixed graph this invariant is a finite bounded formula and hence Booleanizable.

\begin{theorem}[Deterministic contract-DAG composition]\label{thm:dag}
Let \(D\) be a fixed finite acyclic contract graph.  Assume:

\begin{enumerate}
  \item every demand and resource function is primitive recursive;
  \item every local validity predicate is bounded;
  \item every node carries an actual \(\NPrf\) or \(\Prf\,\PRAAx\) proof of its local contract;
  \item the external inputs satisfy the root preconditions.
\end{enumerate}

Then one can effectively construct a fixed primitive-recursive global program \(F_D\) and an object derivation of the global postcondition.  If the graph specification is presented at compile time, the proof compiler is a finite structural recursion and does not require an object-language universal evaluator.
\end{theorem}

\begin{proof}
Compute the reverse-demand program by primitive recursion on the reversed topological list.  The maximum update is primitive recursive.  The demand-adequacy equations follow by a finite induction on that list; for a compile-time fixed graph, the proof generator may simply compose the local equations finitely.

Construct the forward evaluator by primitive recursion on the node index.  At step \(i\), decode predecessor values from the table, apply the fixed term for \(f_i\), and store the result.  Booleanize the global invariant described above.  The base proof follows from table initialization and root assumptions.  For the step, use the demand-adequacy equation to instantiate the local contract of node \(i\); combine its object proof with the induction hypothesis; update the table; and prove that all older entries remain unchanged.  These are equations of fixed primitive-recursive operations and bounded tests.

Apply quantifier-free induction to obtain the invariant after \(m\) steps.  Project the terminal postcondition.  Every step of this argument can be represented by the ordinary proof constructors and the existing primitive-recursion equations.  Equational subproofs may instead be generated in \(\NPrf\) and transferred.  Therefore the compiler returns an object derivation, not merely a host theorem about the evaluator.
\end{proof}

\subsection{Clocked traces for input-coded computations}

A graph or computation may itself vary with the input.  OCiPRA does not attempt to evaluate arbitrary code without a clock.  It uses a fixed transition function
\[
  s_{t+1}=\mathsf{step}(x,s_t)
\]
with fixed initialization, fixed acceptance test, and primitive-recursive bound \(T(x)\).  A trace is a list of \(T(x)+1\) states.

\begin{theorem}[Clocked trace reflection]\label{thm:trace}
Suppose \(\mathsf{init}\), \(\mathsf{step}\), \(\mathsf{inv}\), \(\mathsf{accept}\), and \(T\) are fixed primitive-recursive symbols.  Suppose object proofs establish:
\[
\begin{aligned}
 P(x)&\to \mathsf{inv}(x,\mathsf{init}(x)),\\
 \mathsf{inv}(x,s)&\to \mathsf{inv}(x,\mathsf{step}(x,s)),\\
 \mathsf{inv}(x,s_{T(x)})&\to Q(x,s_{T(x)}).
\end{aligned}
\]
Then there is an effectively constructed \(\Prf\,\PRAAx\) derivation of
\(P(x)\to Q(x,s_{T(x)})\).
\end{theorem}

\begin{proof}
Define the state iterator by primitive recursion.  Booleanize \(\mathsf{inv}\) if it is not already an equality of a Boolean test.  Quantifier-free induction proves the invariant for every counter \(t\).  Universal elimination at \(T(x)\) and the finish proof give \(Q\).  The clock is essential: it makes the state iterator a fixed primitive-recursive function.  Nothing in the construction evaluates a purported program until it happens to halt.
\end{proof}

\subsection{Structural composition without a global loop}

For a small fixed DAG, the global loop of \cref{thm:dag} is optional.  One may compose proof-carrying contracts directly.  If
\[
 P\vdash Q(f(x),x)
 \quad\text{and}\quad
 Q(y,x)\vdash R(g(y,x),y,x),
\]
substitute \(f(x)\) into the second proof, cut with the first, and package the composed program \(g(f(x),x)\).  This route yields shorter proof terms.  The clocked loop is preferable when the number of nodes is input-dependent or when one wants a uniform trace certificate.

\subsection{Master internalization theorem}

\begin{theorem}[OCiPRA internalization theorem]\label{thm:master-compiler}
Consider a mathematical construction normalized into a finite contract specification with the following data:

\begin{enumerate}
  \item fixed primitive-recursive encodings for all finite objects;
  \item either a fixed contract DAG or a fixed clocked transition system;
  \item primitive-recursive demand and resource bounds;
  \item bounded validity, compatibility, and acceptance predicates;
  \item local object derivations for computational equations and logical contracts;
  \item for relational outputs, either a bounded viability proof or a supplied witness certificate;
  \item for representative-independent output, an exact stability proof or quantitative diameter bound;
  \item every nonuniform producer datum exposed as an explicit hypothesis or certificate field.
\end{enumerate}

Then a proof-producing host compiler can construct an \(\NPrf\) or
\(\Prf\,\PRAAx\) derivation of the normalized global theorem.  The construction uses no countable choice, no unbounded compactness, no universal unclocked evaluator, and no reflection from semantic truth.
\end{theorem}

\begin{proof}
Apply \cref{lem:coding} to all finite states.  Apply \cref{thm:booleanization} to every bounded verifier and invariant.  Deterministic portions compile by \cref{thm:dag} or direct structural composition; input-dependent finite computations compile by \cref{thm:trace}.  Relational portions compile by the viability theorem proved in \cref{thm:viability}; a least viable successor is obtained from \cref{thm:least-search}.  Exact or approximate erasure of internal choices is justified only by \cref{thm:stable-exact} or \cref{thm:stable-approx}.  The remaining producer assumptions are present in the precondition, so implication introduction yields the advertised conditional theorem.  Equational proof fragments generated in the native backend transfer through \code{native\_iff\_pra}; all other fragments are already \(\Prf\,\PRAAx\).  Finite proof composition completes the derivation.
\end{proof}

\begin{warning}[No semantic shortcut]
The premise ``the Python checker accepts'' is not a local object derivation.  The compiler theorem becomes applicable only after acceptance and each transition invariant have object-language proofs.  This is the principal remaining engineering task for many repository packets.
\end{warning}

\section{Demand closure, changing parameters, and finite resource bounds}

The reverse-demand pass of a directed acyclic graph is immediate because each edge is visited once.  More difficult presentations contain feedback at the level of \emph{precision requirements}: asking for a more accurate output at one interface may increase the precision needed from another, which in turn raises the first demand.  This is not computational nontermination if the demands are known to lie in a finite box.  The appropriate construction is a finite monotone closure, not an appeal to an abstract least-fixed-point theorem.

\subsection{Monotone closure in a finite box}

Fix a dimension $r>0$, a target vector $a\in\Nat^r$, a certified upper vector $H\in\Nat^r$ with $a\leq H$ componentwise, and a primitive-recursive map
\[
  F:\Nat^r\longrightarrow\Nat^r.
\]
The intended reading is that $d$ is the current demand vector and $F(d)$ is the additional demand forced by all local contracts.  Define the inflationary update
\[
  T(d)=\max(a,d,F(d)),
\]
where maximum is componentwise.  Starting at $d_0=a$, put $d_{n+1}=T(d_n)$.

\begin{theorem}[Finite demand closure]\label{thm:demand-closure}
Assume that $F$ is monotone on the finite box $[0,H]$ and that every iterate $d_n$ generated before stabilization lies below $H$.  Then the sequence stabilizes after at most
\[
  1+\sum_{i<r}(H_i-a_i)
\]
updates.  Its terminal value $d^*$ is the least vector in $[a,H]$ satisfying $F(d^*)\leq d^*$.  The sequence, terminal value, and replay certificate are primitive recursive, and the proof is formalizable by quantifier-free induction in iPRA.
\end{theorem}

\begin{proof}
The update is inflationary, so $d_n\leq d_{n+1}$.  If $d_{n+1}\neq d_n$, at least one coordinate increases by at least one.  Hence the rank
\[
  R(d)=\sum_{i<r}(d_i-a_i)
\]
strictly increases at every nonstationary step and is bounded by
$\sum_i(H_i-a_i)$.  Consequently a stationary point occurs within the stated number of updates.  At stationarity,
$d^*=\max(a,d^*,F(d^*))$, so $a\leq d^*$ and $F(d^*)\leq d^*$.

For minimality, let $e$ satisfy $a\leq e\leq H$ and $F(e)\leq e$.  We show by induction that $d_n\leq e$.  The base case is $a\leq e$.  If $d_n\leq e$, monotonicity gives $F(d_n)\leq F(e)\leq e$, and therefore
$T(d_n)=\max(a,d_n,F(d_n))\leq e$.  Thus $d^*\leq e$.

All vectors have externally fixed length $r$ and may be coded by nested pairing.  Componentwise comparison, maximum, sum, $F$, and iteration for the displayed fuel are fixed primitive-recursive programs.  The invariant
\[
 d_n\leq H\ \wedge\ a\leq d_n\ \wedge\
 R(d_n)\leq \sum_i(H_i-a_i)
\]
is bounded.  Booleanize it and use quantifier-free induction.  A second bounded induction proves minimality relative to a coded candidate $e$.  No impredicative lattice construction occurs.
\end{proof}

The hypothesis that iterates remain below $H$ is itself a useful certificate boundary.  A producer can supply $H$ together with a bounded proof that $T$ preserves the box.  If an update exceeds $H$, the checker rejects the certificate rather than silently increasing its resources.

\begin{remark}[Monotonicity need not be guessed]
For a fixed small box, monotonicity can be checked exhaustively.  In serious applications it is preferable to prove monotonicity compositionally from the demand transformers of the nodes.  The Python companion offers exhaustive checking only as a diagnostic.
\end{remark}

\subsection{Precision-demand composition}

Suppose a node $f$ consumes observations of $X_1,\ldots,X_k$ and produces an observation of $Y$.  Its demand transformer is
\[
  \delta_f(p)=(p_1,\ldots,p_k).
\]
If $g$ consumes the output of $f$ at demand $q=\delta_g(p)$, then the composite demand is
$\delta_f(q)$.  If the same input is used at several sites, its global demand is the maximum of all incoming requirements.  Thus an acyclic network computes demands by reverse composition and maximum.  A cyclic specification computes the least closed demand by \cref{thm:demand-closure}, provided a finite upper box is certified.

A resource transformer behaves similarly, except that resources add or take a maximum depending on whether subcomputations are sequential or reusable.  OCiPRA keeps the two effects separate:
\[
  f:X\xRightarrow[\rho_f]{\delta_f}Y.
\]
A precision guarantee without fuel does not define a primitive-recursive verifier, while fuel without sufficient input precision does not establish the mathematical postcondition.

\subsection{Changing-parameter recursion}

Several repository packets encounter recurrences in which the parameter needed by the preceding stage changes with the stage.  The schematic form is
\begin{equation}\label{eq:changing-rec}
 X(0,k)=b(k),\qquad
 X(n+1,k)=G\bigl(n,k,X(n,D(n,k))\bigr),
\end{equation}
where $b,D,G$ are fixed primitive-recursive functions.  This is not directly the defining equation of ordinary primitive recursion in the displayed parameter $k$, because the recursive call is made at $D(n,k)$ rather than $k$.  A backward-demand/forward-replay compilation removes the difficulty.

For a requested pair $(n,k)$, define a finite schedule
\[
 K_n=k,\qquad K_i=D(i,K_{i+1})\quad(i<n).
\]
Then compute forward
\[
 V_0=b(K_0),\qquad
 V_{i+1}=G(i,K_{i+1},V_i).
\]
The output is $V_n$.

\begin{theorem}[Changing-parameter recurrence compiler]\label{thm:changing-rec}
For fixed primitive-recursive $b,D,G$, the backward schedule and forward replay above define a primitive-recursive function $\widehat X(n,k)$.  iPRA proves both equations in \eqref{eq:changing-rec}.  Moreover, any local invariant preserved by $G$ along the scheduled calls is preserved by $\widehat X$ through an object-language quantifier-free induction.
\end{theorem}

\begin{proof}
The schedule can be stored as a bounded list.  Starting with the singleton list $[k]$, primitive recursion for $n$ steps prepends the value obtained by applying $D$ at the appropriate reversed stage.  Alternatively, encode the whole schedule by a finite table and fill positions from $n$ down to $0$.  Since the loop length is exactly $n$, the schedule function is primitive recursive.

The forward replay is ordinary primitive recursion on $i\leq n$: read $K_i$ and $K_{i+1}$ from the table and apply $G$.  Thus $\widehat X$ is primitive recursive by composition.

At $n=0$, the schedule contains $K_0=k$, so the forward base equation gives $\widehat X(0,k)=b(k)$.  For the successor equation, the schedule for $(n+1,k)$ has $K_{n+1}=k$ and $K_n=D(n,k)$.  Its prefix through stage $n$ is exactly the schedule used for $(n,D(n,k))$.  The forward final step therefore gives
\[
 \widehat X(n+1,k)
 =G\bigl(n,k,\widehat X(n,D(n,k))\bigr).
\]
The object proof expands the two fixed primitive-recursive table programs, proves the prefix identity by quantifier-free induction, and transports the final equation.  For an invariant $I(i,K_i,V_i)$, prove the base instance from the contract for $b$ and the step from the contract for $G$; Booleanization followed by quantifier-free induction yields the result for every $i\leq n$.
\end{proof}

This theorem subsumes many ``adaptive precision'' recurrences.  The backward pass computes which input precision each stage will eventually need; the forward pass performs the actual arithmetic only after all demands are known.  It is also the right implementation model when a naive forward algorithm would repeatedly restart earlier stages at higher precision.

\subsection{Variable-length loops}

A fixed step program $S(x,i,s)$, initial state $s_0(x)$, and primitive-recursive length $L(x)$ define
\[
 s(x,0)=s_0(x),\qquad s(x,i+1)=S(x,i,s(x,i)).
\]
Because $L(x)$ is a term, quantifier-free induction proves an invariant for all $i$ and universal elimination instantiates it at $L(x)$.  The loop length may vary with the input; what matters is that it is a fixed primitive-recursive function.  This observation covers bounded scans, finite-grid methods, exact rational arithmetic routines, Jacobi sweeps with explicit fuel, finite elimination matrices, and certificate replay.

\subsection{Controlled descent and the fixed-dimension boundary}

A termination proof may use a rank in $\Nat^d$ with lexicographic decrease and a primitive-recursive envelope on each coordinate.  For every externally fixed $d$, a primitive-recursive fuel bound can be obtained by nested iteration: the last coordinate supplies an inner counter, the preceding coordinate controls how often that counter can reset, and so on.  The proof is a finite induction on the fixed dimension in the metatheory, producing one fixed primitive-recursive term.

The qualifier ``fixed'' is essential.  An arbitrary input-dependent dimension can encode recursions beyond primitive recursion.  Therefore OCiPRA accepts one of the following:

\begin{itemize}
  \item a fixed rank dimension and proved coordinate envelopes;
  \item a directly supplied primitive-recursive fuel bound;
  \item a finite trace certificate whose length is separately bounded.
\end{itemize}

It rejects the informal assertion ``the lexicographic order is well founded, so the algorithm is primitive recursive.''  Well-foundedness establishes termination in a stronger metatheory; it does not by itself produce the primitive-recursive resource majorant required by iPRA.

\subsection{Formalization obligations}

A generic Lean implementation of \cref{thm:demand-closure,thm:changing-rec} needs four layers:

\begin{enumerate}
  \item fixed-length vector coding and object proofs for projections, update, maximum, and rank;
  \item exact primitive-recursive syntax for the closure and schedule programs;
  \item Booleanized box-preservation and monotonicity contracts;
  \item generated \code{Prf PRAAx} proofs of stabilization, leastness, schedule-prefix equality, and recurrence equations.
\end{enumerate}

The existing repository already reports checked pieces of changing-parameter compilation and backward-demand equations.  The remaining value of a common OCiPRA module is to give those pieces one reusable interface and to connect them to relational viability and stable observables.

\section{Finite relational computation and stable observables}

Many constructive obstructions arise only because a classical proof insists on naming one representative.  Polynomial roots collide, eigenspaces rotate, chart choices overlap, local lifts are nonunique, and compactness arguments yield finite covers rather than a preferred point.  OCiPRA treats the finite relation as primary and eliminates the choice only when the final observable is certified to be stable.

\subsection{Finite cloud systems}

Fix a horizon $H$.  At stage $i\leq H$ there is a nonempty finite cloud
\[
 C_i=\{c_{i,0},\ldots,c_{i,B_i-1}\}.
\]
The candidates and bounds are produced by fixed primitive-recursive functions.  For $i<H$, let
$E_i(a,b)$ be a bounded compatibility predicate between a candidate at stage $i$ and one at stage $i+1$.  Let $G(c)$ be a bounded terminal-good predicate.

A successful transcript is a sequence
$(a_0,\ldots,a_H)$ with $a_i<B_i$, all edges
$E_i(c_{i,a_i},c_{i+1,a_{i+1}})$, and $G(c_{H,a_H})$.
No infinite tree and no countable branch are involved.

Define viability backwards:
\[
\begin{aligned}
 V_H(a)&\;\Longleftrightarrow\;a<B_H\wedge G(c_{H,a}),\\
 V_i(a)&\;\Longleftrightarrow\;a<B_i\wedge
   \exists b<B_{i+1}\bigl(E_i(c_{i,a},c_{i+1,b})\wedge V_{i+1}(b)\bigr).
\end{aligned}
\]
Each $V_i$ is bounded and therefore Booleanizable.

\begin{theorem}[Finite backward viability]\label{thm:viability}
For every fixed primitive-recursive finite cloud system, there are primitive-recursive Boolean programs computing all $V_i$.  iPRA proves:

\begin{enumerate}
  \item $V_i(a)$ holds exactly when $c_{i,a}$ begins a successful suffix to horizon $H$;
  \item if some $V_0(a)$ holds, bounded least search constructs a least successful transcript;
  \item replay of that transcript verifies every edge and the terminal predicate.
\end{enumerate}
\end{theorem}

\begin{proof}
At the terminal layer, viability is the Boolean conjunction of the index bound and $G$.  Given a Boolean program for $V_{i+1}$, compute $V_i(a)$ by scanning the bounded range $b<B_{i+1}$ and taking the Boolean disjunction of compatibility and successor viability.  This is a fixed primitive-recursive fold.

Prove the suffix characterization by reverse induction on the finite distance $H-i$.  At distance zero it is exactly the definition of terminal viability.  In the successor case, $V_i(a)$ supplies a compatible viable successor, and the induction hypothesis supplies its successful suffix; concatenation gives a successful suffix from $a$.  Conversely, the first edge of any successful suffix points to a successor whose remaining suffix is successful; the induction hypothesis makes that successor viable, so the bounded existential defining $V_i(a)$ holds.  Reverse induction can be implemented as ordinary induction on $j=H-i$ using truncated subtraction and the coded viability table.

Assume an initial viable index exists.  Choose the least one by \cref{thm:least-search}.  At stage $i$, scan the finite next layer and choose the least index satisfying both compatibility with the current candidate and $V_{i+1}$.  The defining equation of viability proves that this scan succeeds.  Primitive recursion for $H$ stages builds the transcript.  The maintained invariant states that the current index is viable and that the stored prefix has valid edges.  It is bounded and hence admits quantifier-free induction.  At the last stage, viability is terminal goodness.  This also proves replay correctness.
\end{proof}

\begin{remark}[Finite horizon, not dependent choice]
The least successor at stage $i$ is chosen from a bounded nonempty set whose nonemptiness is proved by the already computed viability bit.  The construction performs $H$ bounded choices by primitive recursion.  It does not use countable choice or assert an infinite coherent branch.
\end{remark}

\subsection{Exact branch elimination}

A successful transcript may be an implementation detail.  Let $O$ be a primitive-recursive observation of its terminal candidate or of the whole finite transcript.

\begin{theorem}[Exact stable observable]\label{thm:stable-exact}
Assume iPRA proves that any two successful transcripts $\tau,\sigma$ satisfy
\[
 O(\tau)=O(\sigma).
\]
If an initial viable candidate exists, the least viable transcript constructed by \cref{thm:viability} yields a deterministic primitive-recursive output $o$, and iPRA proves that every successful transcript has observable $o$.
\end{theorem}

\begin{proof}
Let $\tau_*$ be the least viable transcript.  Its success is proved by \cref{thm:viability}.  Define $o=O(\tau_*)$.  For an arbitrary successful transcript $\sigma$, instantiate the invariance proof with $\tau_*$ and $\sigma$.  Hence $O(\sigma)=o$.  The construction is deterministic because all bounded searches use the least witness, but correctness depends on invariance, not on the arbitrary order of candidate codes.
\end{proof}

Examples of exact stable observables include a sorted list of eigenvalues despite nonunique eigenvectors, a polynomial factor as an unordered multiset of root clusters, a quotient class, an intersection number, a parity, a finite cardinality, and a Boolean decision separated from the branch boundary.

\subsection{Quantitative branch elimination}

Exact equality is often stronger than needed.  Suppose $O(\tau)$ is a rational code and $d$ is a primitive-recursive rational metric.

\begin{theorem}[Diameter-stable observable]\label{thm:stable-approx}
Assume every two successful transcripts satisfy
\[
 d\bigl(O(\tau),O(\sigma)\bigr)\leq \varepsilon(p)
\]
for an explicit primitive-recursive error bound $\varepsilon(p)$.  Then the observable of the least viable transcript is within $\varepsilon(p)$ of the observable of every successful transcript.  If $\varepsilon(p)\leq 2^{-p}$ and the construction is coherent as $p$ increases, these outputs form a regular name.
\end{theorem}

\begin{proof}
The first assertion is immediate by applying the diameter hypothesis to the successful least transcript and an arbitrary successful transcript.  For coherence across precisions, compare the two selected transcripts by transporting both into a common finer cloud or by using a cross-precision compatibility theorem.  The triangle inequality combines the within-cloud diameter and the restriction error.  Predetermine a precision offset large enough that the sum is at most $2^{-p}$.  All estimates are rational inequalities checked by bounded integer arithmetic.
\end{proof}

The theorem is useful even without a global name.  A fixed-precision result such as ``all valid answers lie in an interval of width $2^{-p}$'' is often the strongest constructively valid statement and is sufficient for downstream finite computations.

\subsection{Isolation gaps and exact selectors}

Suppose a finite cloud approximates several objects.  An exact label can be recovered when a gap separates the desired class.  Let $P(c)$ be a bounded classification predicate and suppose a supplied rational $\gamma>0$ proves that candidates from different classes are at distance at least $\gamma$.  If the cloud diameter is below $\gamma/3$, then all candidates compatible with one seed lie in the same class, and bounded least search gives an exact class label.  This is the generic form of:

\begin{itemize}
  \item selecting a reciprocal only after an apartness bound from zero;
  \item selecting a spectral projector only after a spectral gap;
  \item ordering roots only after a separation bound;
  \item selecting a rotation axis only away from the identity or half-turn singularity;
  \item choosing a chart only after a positive margin inside its domain.
\end{itemize}

Without a gap, OCiPRA returns the cloud, quotient, symmetric invariant, or approximate observable.  It does not derive an exact discontinuous selector from a convergent approximation process.

\subsection{Adaptive relational composition}

A consumer may query a producer repeatedly, with later queries depending on earlier answers.  Fix a finite clock $L(x,p)$.  A consumer state contains the finite transcript so far.  At each step it either halts with an output or issues a query $q$; the producer returns one of finitely many answers satisfying a bounded relation.  Form the layered graph of all locally valid interaction states up to the clock and mark accepting terminal states.  Backward viability computes exactly those states from which some accepting continuation exists.

\begin{theorem}[Finite-use relational composition]\label{thm:relational-composition}
For a fixed clocked consumer and bounded finite producer relation, the set of successful interaction transcripts is primitive recursively enumerable within an explicit bound.  If the initial state is viable, iPRA proves existence of a least successful transcript.  If the final observable is exact- or diameter-stable across all successful transcripts, the internal producer choices can be erased according to \cref{thm:stable-exact,thm:stable-approx}.
\end{theorem}

\begin{proof}
Code the consumer state, query, answer, and step counter as one natural.  The candidate layer at time $i$ is the bounded set of all states whose trace length is $i$ and whose local replay checks pass.  Compatibility is one consumer transition plus one valid producer answer.  Acceptance is the terminal consumer predicate.  Apply \cref{thm:viability}; then apply the relevant stability theorem.  Because the entire interaction is bounded by $L$, all quantifiers are finite.
\end{proof}

This theorem explains how local nonuniqueness can coexist with deterministic mathematics.  What must be deterministic is the claimed observable, not every hidden certificate choice.

\subsection{Shrinking clouds without an infinite branch}

At precision $p$, choose the least candidate that lies on a successful path to a finite horizon $h(p)$.  Suppose:

\begin{enumerate}
  \item every full viable cloud at precision $p$ has diameter at most $2^{-(p+c)}$;
  \item compatible candidates at adjacent precisions move by at most a summable rational budget;
  \item the horizon and all bounds are primitive recursive.
\end{enumerate}

Then the selected candidates form a regular name after choosing a fixed offset $c$.  The proof compares the two finite selected paths by extending both to a common finite horizon and summing the motion budgets.  At no point is one branch selected through all precisions simultaneously.  This finite-horizon formulation is the correct replacement for an informal appeal to Kőnig's lemma or dependent choice.

\subsection{Quotients and presentation independence}

A represented mathematical object often has many finite presentations.  Let $E(a,b)$ be a bounded equivalence or approximation relation.  A function on presentations descends to the quotient when an object proof establishes
\[
  E(a,b)\longrightarrow E'(f(a),f(b)).
\]
For exact observables, $E'$ may be equality; for approximate observables, it is a precision-indexed proximity relation.  The contract compiler composes these extensionality proofs just as it composes ordinary postconditions.  Thus ``presentation independence'' becomes a local proof obligation rather than a global set-theoretic quotient construction.

\subsection{Lean target shape}

The central Lean data should contain bounded candidate and compatibility tests as fixed \code{PR} symbols, not host predicates alone.  A useful target is schematically:

\begin{lstlisting}[language=Lean]
structure ViabilityPackage (n : Nat) where
  horizon candidateBound candidate edge terminal : PR n
  viabilityProgram : PR n
  viability_correct : [] |-[PRAAx] viabilitySpec ...
  leastPathProgram : PR n
  leastPath_correct : [initialViable] |-[PRAAx] successfulPath ...

structure StableObservablePackage (n : Nat) where
  relational : ViabilityPackage n
  observable : PR (...)
  invariant : [successful tau, successful sigma] |-[PRAAx]
    observe tau == observe sigma
\end{lstlisting}

The exact de Bruijn indices and arities must be supplied in the implementation.  The point of the structure is that no instance can be manufactured without the object derivations in the proof fields.

\section{Represented mathematics through finite observations}

The statement calculus is useful only if it supports ordinary mathematical constructions.  This section develops a uniform treatment of real and complex quantities, formal series, measures, finite-dimensional geometry, and compact spaces.  The guiding rule is always the same: expose the finite observation actually consumed by the proof and prove its coherence with finer observations.

\subsection{Regular rational names}

A convenient real-number presentation is a primitive-recursive sequence of rationals $a(p)$ satisfying
\[
 |a(p+1)-a(p)|\leq 2^{-(p+2)}.
\]
The rational is encoded by a signed numerator and positive denominator.  Summing the tail gives
\[
 |a(q)-a(p)|\leq 2^{-(p+1)}\qquad(q\geq p).
\]
All inequalities become integer inequalities after multiplying by positive denominators.  Thus regularity at a fixed pair $(p,q)$ is bounded and decidable.  A theorem about the represented real is stated by quantifying $p$ last and proving the finite rational assertion at that precision.

Two names $a,b$ represent the same real when
\[
 |a(p+c)-b(p+c)|\leq 2^{-p}
\]
for a fixed offset $c$.  An operation is extensional when it maps equivalent input names to equivalent output names.  The operation itself is a fixed primitive-recursive program on rational codes.

\subsection{Arithmetic and precision budgets}

For addition, output precision $p$ is obtained by reading both inputs at $p+2$ and adding their rational approximations.  The two input errors sum to at most $2^{-(p+1)}$, leaving a margin for normalization.  Negation has demand $p$.  For multiplication on a supplied bounded region $|x|,|y|\leq M$, use approximations at precision
\[
 p+c+\lceil\log_2(M+1)\rceil
\]
so that
\[
 |xy-x'y'|
 \leq |x|\,|y-y'|+|y'|\,|x-x'|
\]
fits the output budget.  The bound $M$ is a producer datum or a previously proved invariant.

Reciprocal is not uniformly continuous near zero.  It therefore has a guarded contract:
\[
 |x|\geq 2^{-g}
 \quad\Longrightarrow\quad
 \text{a regular name for }1/x.
\]
Read $x$ sufficiently accurately to preserve a lower bound, invert the nonzero rational approximation, and propagate the error using
\[
 \left|\frac1x-\frac1r\right|
 =\frac{|x-r|}{|x|\,|r|}.
\]
The gap $g$ determines the precision offset.  This is a canonical example of an exact operation becoming available only under an explicit separation certificate.

\begin{proposition}[Finite arithmetic DAG]\label{prop:real-dag}
Every fixed expression built from regular names by addition, subtraction, bounded multiplication, and apartness-guarded reciprocal has a primitive-recursive reverse precision transformer and a finite rational evaluator.  Given object proofs of the input name laws and guards, OCiPRA compiles an iPRA proof of the output regularity and extensionality laws.
\end{proposition}

\begin{proof}
Assign to each syntax node its elementary demand transformer and rational error inequality.  Reverse composition gives all leaf precisions; shared subexpressions take the maximum demand.  Evaluate forward with exact rational arithmetic.  The global error proof is a contract DAG whose local postconditions are the displayed rational inequalities.  Since the expression is fixed, the DAG compiler of \cref{thm:dag} applies.  Regularity compares two evaluations at adjacent output precisions after raising both to a common leaf demand.
\end{proof}

\subsection{Finite sums, products, and formal series}

A finite sum or product with primitive-recursive length is a bounded fold.  A formal power series is represented by finite coefficient prefixes.  An operation states exactly how many input coefficients are required:

\begin{itemize}
  \item the coefficient of degree $n$ in a product uses coefficients through $n$;
  \item differentiation through degree $n$ reads the input through degree $n+1$;
  \item composition through degree $n$ reads only finite prefixes whose bound is computed from the valuation and truncation order;
  \item a triangular recurrence computes coefficient $n+1$ from a finite prefix through $n$.
\end{itemize}

Consequently formal identities are particularly well suited to iPRA: at each degree they reduce to finite polynomial arithmetic.  Analytic convergence is a separate contract, requiring radius and tail data.  Keeping formal and analytic specifications separate prevents a common error in which a formal coefficient identity is treated as a statement about values of an infinite series.

\subsection{Functions and finite grids}

A represented function on a rational box may be given by:

\begin{enumerate}
  \item a primitive-recursive rational evaluator $f(q,p)$;
  \item a uniform approximation guarantee at rational points;
  \item a modulus of uniform continuity $\omega(p)$;
  \item a range bound $M$.
\end{enumerate}

A finite grid of mesh $2^{-\omega(p+c)}$ then determines the function to precision $2^{-p}$.  Supremum bounds, positivity, integration, and finite-cover claims can be phrased as certificates on that grid plus the modulus margin.

For integration over $[0,1]$, choose a rational Riemann sum at a mesh prescribed by the modulus.  The error between upper and lower cell values is bounded by the continuity modulus, and the sum uses exact rational arithmetic.  The output at precision $p$ is therefore primitive recursive.  A theorem about an antiderivative or a differential equation needs additional finite contracts for derivatives and truncation errors; it does not follow from the grid evaluator alone.

\subsection{Roots, factors, and singular branches}

For a polynomial with represented coefficients, exact continuous root labels generally do not exist through collisions.  OCiPRA uses one of three interfaces:

\begin{enumerate}
  \item a finite collection of rational disks whose union contains all roots, with multiplicity data and a residual/Rouché-style verifier;
  \item a symmetric observable, such as coefficients, power sums, or a sorted real spectrum;
  \item an isolated root name guarded by a separation disk or derivative lower bound.
\end{enumerate}

Refinement maps each coarse cluster to compatible finer clusters.  Backward viability resolves local ambiguity over a finite horizon; diameter stability yields approximate roots or factors.  A supplied simple-root certificate gives a contraction or Newton/Hensel contract and hence a single branch.  At a singular point the honest output remains a finite set of charts or a quotient-level object unless a producer supplies additional ramification and separation data.

This formulation resolves the statement-level obstruction in Puiseux, algebraic-function, and eigenvalue packets.  It does not by itself prove the difficult producer theorem that finds all singular charts with the required uniform bounds.

\subsection{Finite rational probability laws}

A finite law is a list of rational outcomes and nonnegative rational weights summing to one.  Expectation, variance, convolution, characteristic-function samples at rational arguments, and tail probabilities at rational thresholds are finite sums.  A represented probability distribution is observed through finite quantization and tail certificates.

For central-limit statements, a robust contract includes:

\begin{itemize}
  \item a finite rational approximation of each summand or of its law;
  \item moment bounds and a truncation/tail modulus;
  \item a quantization error budget;
  \item an anti-concentration or continuity margin for the requested rational threshold;
  \item a primitive-recursive sample size or error majorant.
\end{itemize}

The verifier proves a fixed-precision distributional inequality between finite rational laws.  It avoids treating an arbitrary higher-type measure or a discontinuous exact CDF test as primitive data.  Universal quantification over requested precision is introduced only after the finite proof is obtained.

\subsection{Finite nets and topology}

Abstract compactness is replaced by a finite-net presentation.  A compact metric object comes with a primitive-recursive net $N(p)$, a cover proof at radius $2^{-p}$, and finite comparison data.  A continuous map comes with a modulus converting output precision to input-net precision.  Composition is reverse demand composition.

Many topological arguments have a combinatorial finite core:

\begin{itemize}
  \item polygonal or cubical approximations of curves and surfaces;
  \item parity or degree on a finite grid;
  \item nerve complexes of finite covers;
  \item subdivision with an explicit mesh bound;
  \item finite obstruction certificates for intersections or separation.
\end{itemize}

The theorem must state the approximation and stability margin that transport the finite invariant back to the represented object.  A qualitative appeal to compactness or uniform continuity without a modulus is not an iPRA implementation.

\subsection{Geometry and projective data}

Normalized vectors and angles are discontinuous at degenerate inputs.  A robust geometric presentation uses homogeneous or projective coordinates and finite residual tests.  For example, a rotation may be represented by a rational quaternion approximation rather than a globally selected unit axis.  The axis becomes exact only under a gap excluding the singular rotations.  Similarly, a line is often better represented by a nonzero homogeneous triple modulo scaling than by a slope with a vertical-line exception.

The relational framework permits several locally valid charts.  Stable incidence, orientation with a nonzero determinant margin, squared distances, and polynomial identities descend without choosing one chart globally.

\subsection{Diophantine objects}

For a polynomial system with integer or rational coefficients, evaluation at a bounded candidate is primitive recursive.  Thus a supplied height bound $H$ yields a complete finite search.  More efficient certificates---resultants, norm matrices, annihilator identities, descent data, local solubility witnesses---may reduce the search box, but they remain finite verifiers.

The framework therefore separates two statements:

\begin{enumerate}
  \item \emph{verification theorem}: from a height, cover, or algebra certificate, compute and prove a complete list;
  \item \emph{producer theorem}: prove that every input in a class has such a certificate and compute one with a primitive-recursive bound.
\end{enumerate}

The first is often accessible in iPRA.  The second may contain the entire depth of Faltings, Siegel, Mordell--Weil, or analytic number theory and must not be hidden inside the word ``finite.''

\subsection{Observation towers}

A useful general abstraction is a tower
\[
 X_0 \xleftarrow{r_0} X_1 \xleftarrow{r_1} X_2\xleftarrow{}\cdots
\]
of finite observation sets with restriction maps.  A deterministic name chooses compatible elements; a relational name provides finite nonempty clouds with compatibility edges.  An operation consists, for every requested level $p$, of:

\begin{enumerate}
  \item a finite input level demand $\delta(p)$;
  \item a primitive-recursive map or finite relation on observations;
  \item a local validity proof;
  \item a restriction/coherence proof.
\end{enumerate}

This tower language unifies Cauchy names, jets, finite distributions, covers, truncations, root clusters, spectra, and proof certificates.  OCiPRA contributes the compilation discipline that turns the finite-level proofs into an object derivation and records exactly where an external producer remains.

\section{Resolution of the reported difficulty classes}

This section applies the framework to the recurring obstacles documented in the repository.  A ``resolution'' may be one of three kinds:

\begin{enumerate}
  \item an internal theorem schema that removes the obstruction;
  \item a principled reformulation that identifies the strongest constructive statement;
  \item a sharp producer boundary showing what additional mathematics remains.
\end{enumerate}

It is important not to call the third kind a proof of the missing deep theorem.  Exposing the missing producer is nevertheless progress: it prevents a verifier proof from being mistaken for a producer theorem and gives a precise target for further work.

\subsection{Ambient Lean truth is not an iPRA proof}

\paragraph{Difficulty.}
Many packets contain executable Lean definitions and host-level theorems, but the desired endpoint is a term of \code{Prf PRAAx} or \code{NPrf}.

\paragraph{Resolution.}
Every local contract includes its object derivation as a field.  The global compiler is typed to return an object derivation.  Equation-heavy local proofs may be generated in the native calculus and transferred by \code{native\_iff\_pra}; bounded logical composition remains in \code{Prf PRAAx}.  Semantic testing is retained only as regression evidence.

\paragraph{Residual work.}
A generic Booleanizer, trace compiler, and viability compiler must still be implemented against the exact repository syntax.  The architecture removes the logical ambiguity but not the engineering effort.

\subsection{Induction on an all-precision assertion}

\paragraph{Difficulty.}
A Cauchy, series, or approximation invariant quantified over every precision is not quantifier free, while iPRA induction is restricted to quantifier-free formulas.

\paragraph{Resolution.}
Fix the requested precision $p$.  Reverse-compute a finite table of every input precision and intermediate error budget used at $p$.  Booleanize the finite invariant and perform quantifier-free induction on the bounded loop.  Only after obtaining the finite theorem introduce the universal quantifier over $p$.  The proof never inducts on the all-precision assertion itself.

\subsection{No universal unclocked evaluator}

\paragraph{Difficulty.}
One cannot postulate a total primitive-recursive evaluator that executes every coded primitive-recursive program, or every partial search, without a bound.

\paragraph{Resolution.}
There are two accepted routes.  A fixed source expression is compiled in the metatheory to one fixed \code{PR} term and one proof.  An input-coded computation is represented by a fixed transition function plus a primitive-recursive clock and bounded trace verifier.  The clocked reflection theorem \cref{thm:trace} then applies.

This distinction also clarifies self-interpreters and proof checkers.  A bounded proof checker for a fixed syntax is primitive recursive.  Turning acceptance into an arbitrary mathematical conclusion still requires a separately proved soundness theorem inside the relevant metatheory; OCiPRA never treats checker acceptance as arithmetic truth by reflection.

\subsection{Exact geometric selectors and omniscience}

\paragraph{Difficulty.}
A globally exact axis for a rotation, a sign at zero, a normalized direction through a zero vector, or a sharp branch choice at a collision can encode limited principles of omniscience.

\paragraph{Resolution.}
Return projective data, quaternions, finite axis clouds, or stable polynomial observables unconditionally.  Add an apartness or separation gap for an exact selector.  The finite cloud is refined by \cref{thm:viability}; exact branch elimination uses \cref{thm:stable-exact}.  The gap is not a technical nuisance but the constructive information needed to make the selector continuous on the promised domain.

\subsection{Sharp endpoints and one-sided density}

\paragraph{Difficulty.}
Several approximation statements become nonconstructive when they require equality at a closed endpoint, a negative-only witness, or an exact decision at a boundary.

\paragraph{Resolution.}
Use strict rational margins, positive separation moduli, and finite alternatives.  For example, prove that a point lies within an open error interval, or return either a positive witness or a certificate that all bounded candidates fail by a stated margin.  Closed-endpoint consequences may then be recovered under an extra decidable equality or apartness hypothesis.

\subsection{Finiteness versus height and stopping}
\label{sec:height-obstruction}

\paragraph{Difficulty.}
A theorem that a set of rational or integral points is finite, or even has cardinality at most $N$, does not provide a height bound, a terminating enumeration, or a complete list.

\begin{theorem}[Height-to-list conversion]\label{thm:height-list}
Let $R(x,y)$ be a bounded decidable relation on coded rational or integer candidates.  If a primitive-recursive height bound $H(x)$ is supplied and every solution satisfies $|y|\leq H(x)$ in the chosen coding, then a primitive-recursive program returns a complete list of all solutions, and iPRA proves its completeness.
\end{theorem}

\begin{proof}
Enumerate the finite candidate box determined by $H(x)$, evaluate the bounded relation, and append exactly the passing candidates.  The complete bounded-list theorem \cref{cor:bounded-list} proves soundness and completeness.  Rational candidates require a specified numerator/denominator convention and denominator bound; these are part of the height certificate.
\end{proof}

\begin{proposition}[Cardinality does not uniformly give a cutoff]\label{prop:cardinality-cutoff}
There is no uniform primitive-recursive transformation that takes an arbitrary primitive-recursive decidable predicate known to have at most one witness and returns a bound above its possible witness.
\end{proposition}

\begin{proof}
For a machine or partial recursive computation index $e$, define the primitive-recursive predicate
\[
 R_e(t)\quad\Longleftrightarrow\quad
 \text{``the computation $e$ halts for the first time exactly at step $t$''}.
\]
The predicate has at most one witness.  If a uniform primitive-recursive transformer returned a bound $B(e)$ containing that witness whenever it exists, one could decide whether $e$ halts by simulating through $B(e)$.  This contradicts the undecidability of the halting problem.  The same argument applies to a purported complete-list producer from a mere cardinality bound.
\end{proof}

\paragraph{Framework treatment.}
The verifier theorem is conditional on a height bound or complete-list certificate.  General Faltings and Siegel become producer obligations: prove and compute the requisite bound for the stated class.  A special family with an explicit Runge, norm, gap, or resultant certificate may be fully internalized without solving the general producer problem.

\subsection{Mordell--Weil and descent representatives}

\paragraph{Difficulty.}
Finite Selmer data or a finite collection of covering classes does not automatically provide rational points on the coverings, nor generators of the Mordell--Weil group.

\paragraph{Resolution.}
Represent each needed rational covering point by a finite certificate and bounded verifier.  Once representatives are supplied, the descent map, group arithmetic, relation checking, saturation within a bound, and finite enumeration are contract-DAG computations.  The theorem is conditional on the representative package unless a separate producer is proved.

This separates three logically different outcomes:

\begin{enumerate}
  \item a finite obstruction proving that a class has no rational representative;
  \item a checked representative for each surviving class;
  \item a proof that the resulting points generate the required group.
\end{enumerate}

Only the first two are direct finite certificates; the third needs its own height/index argument.

\subsection{Ineffective analytic constants}

\paragraph{Difficulty.}
Prime-distribution theorems may invoke ineffective constants, exceptional zeros, compactness over character families, or asymptotic notation without a computable threshold.

\paragraph{Resolution.}
Factor the theorem into:

\begin{enumerate}
  \item an explicit analytic-input certificate containing all constants, thresholds, zero-free regions, exceptional data, and finite estimates;
  \item a primitive-recursive verifier for that certificate;
  \item a finite derivation that propagates the estimates to the claimed bound.
\end{enumerate}

The propagation theorem is a natural iPRA target.  The uniform producer for the full analytic certificate remains the deep theorem.  For fixed moduli, finite ranges, or explicitly certified zero data, substantial special cases may become unconditional.

\subsection{Higher-type probability and discontinuous distribution tests}

\paragraph{Difficulty.}
An arbitrary measure is not numerical first-order data, and the exact indicator of a threshold is discontinuous at atoms.

\paragraph{Resolution.}
Use finite rational laws, quantization, tail and moment moduli, and rational thresholds with anti-concentration margins.  At fixed precision every expectation and convolution is a finite rational sum.  The final theorem can be a bounded-Lipschitz estimate or a CDF estimate at thresholds whose boundary mass is explicitly controlled.  This yields a constructive central-limit interface without pretending that every measure comes with effective presentation data.

\subsection{Root ordering and collisions}

\paragraph{Difficulty.}
Exact root orderings, branch labels, eigenvectors, and Puiseux branches cannot remain continuous through multiplicity changes.

\paragraph{Resolution.}
Use isolating clusters, multisets, symmetric functions, sorted real eigenvalue lists, or finite compatible charts.  A derivative/separation certificate permits an isolated branch.  Stable observables are compiled by \cref{thm:stable-exact,thm:stable-approx}; no global exact selector is asserted at collisions.

\subsection{Abstract compactness in topology and analysis}

\paragraph{Difficulty.}
A proof may say ``choose a finite subcover'' or ``by compactness choose a uniform scale'' without giving a primitive-recursive net or modulus.

\paragraph{Resolution.}
The object is presented with an explicit finite-net modulus.  A finite cover certificate names the selected cells and verifies coverage with a margin.  Subdivision depth, mesh, and parity/degree calculations are bounded.  If no net modulus can be produced, the theorem remains conditional on one.

\subsection{Producer--verifier--derivation factorization}

The preceding examples are instances of one theorem-design principle.

\begin{theorem}[Producer--verifier--derivation factorization]\label{thm:pvd}
Suppose a classical proof of $\forall x\in C\,\Phi(x)$ uses finite auxiliary data $c$ only through finitely checkable properties.  Assume:

\begin{enumerate}
  \item a primitive-recursive bounded predicate $\mathsf{Cert}(x,c)$ expresses those properties;
  \item a fixed finite construction transforms a valid $c$ into the desired witness or observation;
  \item local object derivations prove the transformation correct.
\end{enumerate}

Then OCiPRA yields an object derivation of
\[
  \forall x\;\bigl(C(x)\wedge\exists c\,\mathsf{Cert}(x,c)\to\Phi(x)\bigr).
\]
If, in addition, iPRA proves a primitive-recursive producer $p$ with
$C(x)\to\mathsf{Cert}(x,p(x))$, then it proves the unconditional theorem
$\forall x(C(x)\to\Phi(x))$.
\end{theorem}

\begin{proof}
The first statement applies bounded Booleanization, the contract compiler, and existential elimination to an arbitrary certificate satisfying $\mathsf{Cert}$.  For the second, instantiate the conditional theorem with $p(x)$ and cut with the producer proof.  The derivation is entirely intuitionistic.  The substantive difference between the conditional and unconditional theorem is exactly the producer proof.
\end{proof}

\subsection{A diagnostic table}

\small
\begin{longtable}{@{}p{0.16\textwidth}p{0.24\textwidth}p{0.34\textwidth}p{0.18\textwidth}@{}}
\toprule
Area & Obstruction & OCiPRA treatment & Status\\
\midrule
\endfirsthead
\toprule
Area & Obstruction & OCiPRA treatment & Status\\
\midrule
\endhead
Logic & Host theorem is not \code{Prf PRAAx} & Proof-carrying contracts; dual native/full backend & Compiler engineering remains\\
Analysis & All-precision invariant is not QF & Fix precision, finite demand table, QF induction, quantify last & Schema resolved\\
Computation & Universal unclocked evaluator & Compile fixed syntax or verify a clocked trace & Resolved boundary\\
Geometry & Exact axis/direction at singularity & Projective or finite-cloud output; exact under gap & Exact unrestricted claim rejected\\
Approximation & Sharp endpoint or one-sided choice & Strict margins and separation moduli & Positive form resolved\\
Diophantine geometry & Finiteness lacks cutoff & Conditional height/list verifier; producer separated & General producer open\\
Elliptic descent & Classes lack rational representatives & Checked representative certificate & Uniform producer open\\
Prime distribution & Ineffective constants & Explicit analytic-input certificate & Propagation internal; producer open\\
Probability & Higher-type measures, CDF boundary & Finite rational laws, tail and anti-concentration moduli & Effective presentations covered\\
Algebra & Root/eigenvector collisions & Clusters, multisets, stable observables, gap selectors & Statement resolved\\
Topology & Abstract compactness & Finite nets, covers, subdivision and parity certificates & Theorem-specific combinatorics remain\\
\bottomrule
\end{longtable}
\normalsize

The table should be read as an implementation map, not as a declaration that every listed classical theorem has been proved in iPRA.

\section{Worked finite examples}

The examples in this section are deliberately small enough to replay exactly.  They exercise the generic algorithms rather than pretending to settle the deep application theorems.  The Python companion computes the same data with exact integers and rational fractions.

\subsection{A deterministic contract DAG}

Take external inputs $x=3$ and $y=4$.  The fixed graph has three nodes:
\[
 s=x+y,\qquad q=s^2,\qquad o=q+x.
\]
Assign conservative precision demands
\[
\begin{aligned}
 \delta_s(p)&=(p+1,p+1),\\
 \delta_q(p)&=(p+2),\\
 \delta_o(p)&=(p+1,p+1),
\end{aligned}
\]
where the two inputs to $o$ are $q$ and $x$.  Request output demand $2$.  The reverse pass yields
\[
 d_o=2,
 \quad d_q=3,
 \quad d_s=5,
 \quad d_x=d_y=6.
\]
The forward values are
\[
 s=7,
 \quad q=49,
 \quad o=52.
\]
The replay certificate stores, for each node, its demand, dependency values, output, and local-validity bit.  The exact program checks that the reverse demands and every local equation recompute to the stored values.

In an object proof, the three equations are proved by the defining equations for addition and multiplication programs.  Structural contract composition is enough; no global induction is required because the graph is fixed and has only three nodes.  The same representation scales to a generated finite DAG, where a single table invariant is more convenient.

\subsection{Least bounded search}

Consider the relation $R(n)\equiv n^2\geq30$ with bound $10$.  The scan tests
\[
 0,1,2,3,4,5,6,7,8,9,10
\]
and obtains the Boolean trace
\[
 0,0,0,0,0,0,1,1,1,1,1.
\]
The least witness is $6$.  The certificate contains the bound, found-bit, witness, and full test vector.  Its verifier recomputes the vector and checks both witness validity and absence of an earlier hit.  The object proof uses exactly the invariant in \cref{thm:least-search}.

This pattern should replace an unspecified phrase such as ``choose the least valid certificate.''  The candidate bound and Booleanized relation are indispensable inputs.

\subsection{Cyclic demand closure}

Let
\[
 F(a,b)=\left(\min(10,b+1),
               \min\left(10,\left\lfloor\frac{a+1}{2}\right\rfloor+2\right)
        \right),
\]
with target $(2,1)$ and box upper bound $(10,10)$.  Iterating
$T(d)=\max((2,1),d,F(d))$ gives
\[
\begin{split}
 (2,1)&\to(2,3)\to(4,3)\to(4,4)\\
      &\to(5,4)\to(5,5)\to(6,5).
\end{split}
\]
The final vector is stationary.  Every strict update increases the coordinate sum, so the generic rank proof supplies fuel independently of this concrete run.  The Python checker also exhaustively verifies monotonicity on the finite box for this example.

The mathematical point is that mutual precision dependencies need not be manually over-approximated when a certified finite box is available.  The least closed demand preserves more information and often reduces downstream computation.

\subsection{A relational cloud for \texorpdfstring{$\sqrt2$}{sqrt(2)}}

Use four finite layers of rational candidates:
\[
\begin{aligned}
 C_0&=\left\{1,2\right\},\\
 C_1&=\left\{\frac32,\frac43\right\},\\
 C_2&=\left\{\frac75,\frac{17}{12}\right\},\\
 C_3&=\left\{\frac{41}{29},\frac{99}{70}\right\}.
\end{aligned}
\]
Candidates in layers $i$ and $i+1$ are compatible when their distance is at most $2^{-(i+1)}$.  A terminal candidate is good when
$|q^2-2|\leq1/100$.  Both candidates are viable in every layer.  The least viable path is
\[
  1,\quad\frac32,\quad\frac75,\quad\frac{41}{29}.
\]
Its terminal residual is
\[
 \left|\left(\frac{41}{29}\right)^2-2\right|=\frac1{841}.
\]
The two viable terminal values differ by
\[
 \frac{99}{70}-\frac{41}{29}=\frac1{2030}.
\]
Thus their exact floor observable is invariant and equals $1$, while the rational endpoint itself is only determined up to diameter $1/2030$.  The example illustrates both branch-elimination theorems: exact invariance for a discrete observable and quantitative stability for an approximate one.

No claim is made that these four layers alone define $\sqrt2$ as a real.  To obtain a name one supplies such clouds at every precision, cross-precision compatibility, and a shrinking diameter bound.

\subsection{Regular-name reciprocal}

Suppose a regular rational name $a(p)$ represents $x$ and a certificate gives $|x|\geq2^{-g}$.  Read $a$ at precision $m=p+2g+c$ for a fixed safety constant $c$.  Regularity makes
$|a(m)-x|\leq2^{-(m+1)}$.  Choose $c$ so this is at most $2^{-(g+2)}$.  Then
\[
 |a(m)|\geq |x|-|a(m)-x|
          \geq 2^{-g}-2^{-(g+2)}
          \geq 2^{-(g+1)}.
\]
Define the output rational as $1/a(m)$.  The error is bounded by
\[
 \left|\frac1x-\frac1{a(m)}\right|
 \leq 2^{2g+1}|x-a(m)|,
\]
which is below $2^{-p}$ by the choice of $m$.  All quantities are rational or dyadic; the checker clears positive denominators.  Without the gap $g$, the precision transformer cannot be uniform.

\subsection{Complete bounded Diophantine enumeration}

Let $F_1,\ldots,F_s\in\mathbb Z[X_1,\ldots,X_n]$, and suppose a certificate proves that every integer solution satisfies $|x_i|\leq H_i$.  Enumerate the box
\[
  \prod_{i=1}^n[-H_i,H_i],
\]
evaluate all polynomials exactly by Horner schemes, and retain precisely the common zeros.  The loop count is
$\prod_i(2H_i+1)$, a primitive-recursive resource bound.  A table invariant proves that after scanning the first $k$ coded candidates, the output list contains exactly the solutions among those candidates.

The verifier theorem is unconditional once the $H_i$ and their coverage proof are supplied.  Producing useful $H_i$ may require Runge estimates, norm identities, height theory, descent, or another deep argument.  OCiPRA records that dependency in the precondition rather than conflating it with the finite scan.

\subsection{Spectral data}

For a rational symmetric matrix, a finite Jacobi or grid certificate may produce an approximate orthogonal matrix $Q$ and diagonal $D$ with a small residual
$\|A-QDQ^T\|$.  Eigenvectors are unstable at repeated eigenvalues, but the sorted diagonal list is stable under an explicit comparison estimate.  Therefore:

\begin{itemize}
  \item the certificate relation may retain many valid $Q$ values;
  \item the stable observable is the sorted eigenvalue list with multiplicity;
  \item a spectral projector onto a cluster is selected only under a gap separating that cluster.
\end{itemize}

The relational contract removes the need to choose a coherent eigenbasis through collisions.  Matrix-specific inequalities still require their own object derivations; this example explains the correct target statement.

\subsection{A conditional Faltings/Siegel interface}

For a curve code $C$, let $\mathsf{HeightCert}(C,c)$ assert that $c$ contains a bound $H$, a proof payload accepted by a fixed bounded checker, and the conclusion that every rational point has numerator and denominator height at most $H$.  Let $L(C,c)$ be the complete finite list obtained by bounded enumeration.

The internal theorem has the form
\[
 \mathsf{Curve}(C)\wedge\mathsf{HeightCert}(C,c)
 \longrightarrow
 \forall P\;\bigl(P\in C(\mathbb Q)\leftrightarrow P\in L(C,c)\bigr).
\]
This is meaningful and useful: it turns any future height producer into a complete decision procedure by one cut.  It is not the general Faltings or Siegel theorem, because existence of $c$ has not been established uniformly.

\subsection{An analytic-input certificate}

A prime-distribution propagation theorem may accept a tuple containing:

\begin{itemize}
  \item a modulus range and numerical threshold;
  \item zero-free-region constants;
  \item an exceptional-character flag and exceptional-zero interval, when relevant;
  \item verified finite bounds for short ranges;
  \item rational upper bounds for every error term used in the proof.
\end{itemize}

The checker verifies the finite inequalities and then runs a contract DAG that combines them.  The result can be an explicit inequality for all $x$ above the supplied threshold, represented in integer/rational arithmetic.  A theorem producing such tuples uniformly is a separate analytic goal.  This formulation permits honest progress on the propagation and bookkeeping layer without claiming an effective constant that has not been derived.

\subsection{Admissibility classifier output}

The executable companion classifies three sample theorem profiles:

\begin{enumerate}
  \item a fixed finite matrix identity with explicit resources and local proofs is \emph{object-internalizable};
  \item a Faltings-style verifier lacking a height producer is \emph{conditional on a producer};
  \item an exact root label through a collision must be \emph{reformulated relationally} or guarded by a separation gap.
\end{enumerate}

The classifier is intentionally conservative and syntactic.  It cannot prove the mathematics.  Its purpose is to force the theorem author to name the missing datum before a long formalization begins.

\section{Lean formalization architecture}

The formalization target is the repository's actual \code{LeanProofs.ConciseIPRA} namespace.  The package does not introduce a surrogate proof relation.  This section gives enough definitions and theorem signatures to make the path to a complete formalization explicit.

\subsection{Existing endpoints to reuse}

The current library supplies:

\begin{itemize}
  \item \code{PR n}, the syntax of fixed $n$-ary primitive-recursive functions;
  \item \code{PTm n}, ordinary iPRA terms;
  \item \code{Fm PTm n}, formulas with $n$ free variables;
  \item \code{PRAAx}, the theory containing defining equations, successor facts, and quantifier-free induction;
  \item the derivability judgment \code{Prf PRAAx}, written with the repository sequent notation;
  \item \code{NFm} and \code{NPrf}, the native scalar formulas and proof relation;
  \item \code{native\_iff\_pra}, transporting native derivability to and from ordinary iPRA;
  \item \code{PRA.qfInd}, the exact object induction rule needed by the finite compilers.
\end{itemize}

The repository also contains examples constructing actual object derivations, including a proof of $\forall x\,(x+0=x)$ and an existential identity contract.  These examples are the right integration tests because their conclusions are sequents, not host-level arithmetic propositions.

\subsection{Proof-carrying contract interfaces}

The supplied file \code{OCiPRA/Contracts.lean} defines the following core structures.

\begin{lstlisting}[language=Lean]
structure Booleanized (n : Nat) where
  formula : Fm PTm n
  test : PR n
  equivalent : [] |-[PRAAx]
    Fm.iff formula (accepts test)

structure ProgramContract (n : Nat) where
  pre : Fm PTm n
  program : PR n
  post : Fm PTm (n + 1)
  certified : [pre] |-[PRAAx]
    post.subst (inst (PTm.app program Syn.var))

structure NativeContract (n : Nat) where
  assumptions : List (NFm n)
  conclusion : NFm n
  proof : NPrf assumptions conclusion
\end{lstlisting}

The theorem
\begin{lstlisting}[language=Lean]
theorem NativeContract.toPRA (c : NativeContract n) :
  c.assumptions.map NFm.toPRA |-[PRAAx]
    c.conclusion.toPRA :=
  native_iff_pra.mp c.proof
\end{lstlisting}
implements the dual-backend bridge.  A deterministic contract obtains its existential consequence by \code{Prf.exI}, and a conditional theorem by \code{Prf.impI}.  A quantifier-free induction package stores the invariant, its \code{QF} proof, and actual base and step derivations; its \code{conclude} theorem calls \code{PRA.qfInd}.

These structures are not axioms.  Their proof fields must be populated before an instance exists.

\subsection{Phase A: bounded formula syntax}

Introduce an inductive syntax that mentions fixed object terms rather than arbitrary host predicates.  One possible shape is:

\begin{lstlisting}[language=Lean]
inductive BFm : Nat -> Type
  | eq  (s t : PTm n) : BFm n
  | and (A B : BFm n) : BFm n
  | or  (A B : BFm n) : BFm n
  | imp (A B : BFm n) : BFm n
  | bot : BFm n
  | bex (bound : PTm n) (body : BFm (n+1)) : BFm n
  | ball (bound : PTm n) (body : BFm (n+1)) : BFm n
\end{lstlisting}

Define \code{BFm.toFm}.  Define \code{BFm.test} by generating a \code{PR} program.  Atomic equality needs a reusable primitive-recursive equality test and object proof.  Connectives use Boolean combinators.  Bounded quantifiers generate folds whose state is a bit.

The main endpoint is:

\begin{lstlisting}[language=Lean]
theorem BFm.booleanize (phi : BFm n) :
  [] |-[PRAAx] Fm.iff phi.toFm (accepts phi.test)
\end{lstlisting}

The recursive connective cases use the intuitionistic proof constructors and finite Boolean truth-table lemmas.  The bounded-quantifier cases instantiate a generic fold-induction lemma.  It is preferable to prove that lemma once rather than reproduce de Bruijn substitution calculations in every quantifier case.

\subsection{Phase B: a library of object-level folds}

Define fixed PR syntax and object equations for:

\begin{itemize}
  \item Boolean conjunction, disjunction, complement under a Boolean guard;
  \item pair-state projection and update;
  \item bounded maximum and sum;
  \item finite list/table lookup and update;
  \item a generic fixed-step fold and its invariant rule.
\end{itemize}

A reusable fold theorem should have an endpoint analogous to:

\begin{lstlisting}[language=Lean]
def compileFold
    (init step : PR ...)
    (inv : Fm PTm ...)
    (qf : inv.QF)
    (base : ... |-[PRAAx] ...)
    (preserve : ... |-[PRAAx] ...) :
    ... |-[PRAAx] Fm.all inv
\end{lstlisting}

The precise arities should be calculated by helper definitions and checked with small examples.  The repository's de Bruijn representation makes scoping errors more likely than mathematical errors, so every helper needs substitution simp lemmas.

\subsection{Phase C: least bounded search}

Define the scan state as a coded pair \code{(found,witness)}.  The relation is supplied as a \code{Booleanized (n+1)} and the bound as a \code{PR n}.  Generate a PR program \code{least}.  The invariant is the Booleanization of:

\begin{align*}
 \mathsf{found}=0&\leftrightarrow
      \forall z<k\;\neg R(z),\\
 \mathsf{found}=1&\to
      \mathsf{witness}<k\wedge R(\mathsf{witness})
      \wedge\forall z<\mathsf{witness}\;\neg R(z).
\end{align*}

The exact endpoint is:

\begin{lstlisting}[language=Lean]
theorem leastBounded_correct
    (R : BFm (n+1)) (B : PR n) :
  [existsWithin R B] |-[PRAAx] leastPost R B
\end{lstlisting}

A second theorem exposes the failure convention with a found-bit, preventing downstream code from treating the default zero as a witness.

\subsection{Phase D: fixed DAG compiler}

The DAG syntax lives in Lean's metatheory.  It is a topologically sorted finite vector of typed nodes.  Each node records:

\begin{itemize}
  \item arity and predecessor indices;
  \item a fixed \code{PR} program;
  \item demand and resource syntax;
  \item a \code{ProgramContract}, \code{NativeContract}, or \code{RelationalContract};
  \item substitution maps identifying local variables with table entries.
\end{itemize}

For a small fixed DAG, define \code{compileDAG} by recursion on the host vector.  At each node, weaken and substitute its local proof, then compose by cuts and equality transport.  The return type is itself the desired sequent:

\begin{lstlisting}[language=Lean]
def compileDAG (D : FixedDAG n)
    (local : D.LocalObjectProofs) :
    [D.pre] |-[PRAAx] D.instantiatedPost
\end{lstlisting}

For an input-coded node list, generate one clocked table evaluator and use the fold compiler.  Do not define an object-level evaluator for arbitrary \code{PR} syntax.

\subsection{Phase E: demand closure and changing parameters}

Encode fixed-length vectors by nested pairs or bounded tables.  Implement
\[
 T(d)=\max(a,d,F(d))
\]
and the rank sum.  The stabilization theorem needs:

\begin{enumerate}
  \item a proof that the initial vector lies in the box;
  \item an object proof that $T$ preserves the box;
  \item an object proof of monotonicity for leastness;
  \item a rank-increase lemma for every nonstationary step.
\end{enumerate}

The changing-parameter compiler should reuse the repository's existing term-to-symbol and backward-demand work.  Add the missing global prefix theorem and the forward replay invariant, then expose the result as a \code{ProgramContract}.

\subsection{Phase F: finite viability}

Define a fixed cloud interface by primitive-recursive candidate bounds, candidate decoding, compatibility, and terminal tests.  The backwards viability table is a bounded Boolean fold.  The main targets are:

\begin{lstlisting}[language=Lean]
theorem viability_iff_successfulSuffix
    (C : FixedCloud n) :
  [] |-[PRAAx] C.viabilitySpecification

theorem viablePath_correct
    (C : FixedCloud n)
    (h : C.LocalObjectProofs) :
  [C.initialViable] |-[PRAAx] C.pathPost
\end{lstlisting}

The second theorem uses the least-bounded-search compiler at each stage.  Its invariant records current viability and validity of the prefix.  A separate replay theorem should permit exported path certificates to be checked without re-running the producer.

\subsection{Phase G: stable observables}

Define exact and quantitative stability packages.  An exact package carries an actual derivation comparing two successful transcript codes.  The compiler instantiates it with the least transcript and an arbitrary transcript.  A quantitative package adds rational metric and triangle-inequality contracts.

\begin{lstlisting}[language=Lean]
theorem stableObservable_correct
    (C : FixedCloud n) (s : C.ExactStable) :
  [C.pre] |-[PRAAx] C.observablePost
\end{lstlisting}

There must be no default instance that turns any relational contract into a deterministic contract.  This negative design property prevents exact-choice claims from entering through typeclass inference or an unnoticed classical choice in host code.

\subsection{Phase H: represented-object adapters}

After the generic compiler exists, add adapters in increasing difficulty:

\begin{enumerate}
  \item fixed finite arithmetic and existing observation combinators;
  \item regular rational names and guarded reciprocal;
  \item finite coefficient prefixes and formal-series operations;
  \item root clusters and stable spectral lists;
  \item finite nets and finite rational laws;
  \item height/list and descent certificate verifiers.
\end{enumerate}

Each adapter should expose its demand transformer, resource bound, object local proofs, and producer fields.  It should not import a large semantic theorem and then label a host proposition as the iPRA result.

\subsection{Proof export and independent replay}

The repository has encoded proof-tree machinery.  Once a generated \code{Prf PRAAx} value exists, a second export layer should:

\begin{enumerate}
  \item encode the formula, context, and proof tree;
  \item write a compact certificate;
  \item replay it with the existing checker;
  \item compare the decoded conclusion with the intended target formula.
\end{enumerate}

This is not needed for logical soundness inside Lean, but it provides a small independently executable proof artifact and detects accidental dependence on host elaboration conveniences.

\subsection{Audit gates}

A theorem is marked complete only when all of the following hold:

\begin{itemize}
  \item the conclusion is an actual \code{Prf PRAAx} or \code{NPrf} judgment;
  \item \code{\#print axioms} reports no \code{sorryAx} and no custom axiom;
  \item every producer assumption occurs visibly in the formula or input record;
  \item every loop has a fixed PR program or explicit primitive-recursive clock;
  \item exact branch elimination cites a stability or isolation derivation;
  \item a small integration example compiles against the root repository workspace;
  \item ideally, the encoded proof tree also replays successfully.
\end{itemize}

The supplied Lean sources contain no \code{axiom} and no \code{sorry}.  They were not compiled in the artifact environment because no Lean executable was available.  That status is deliberately weaker than the status of the mathematical arguments and Python tests.

\section{Limits, negative results, and further research}

A useful foundational framework must say not only what it compiles, but also what it refuses to infer.  The restrictions below are not incidental implementation gaps.  They mark real differences between finite verification, effective production, and discontinuous choice.

\subsection{What the framework does not manufacture}

OCiPRA does not manufacture:

\begin{itemize}
  \item a primitive-recursive clock from a proof that an arbitrary computation terminates;
  \item a height bound from qualitative finiteness or a cardinality bound;
  \item a convergence modulus from bare convergence or metastability without a rate-producing argument;
  \item a finite net from abstract compactness without presentation data;
  \item rational points on descent coverings from their finite isomorphism classes;
  \item effective analytic constants from an ineffective existence proof;
  \item exact root, eigenvector, axis, sign, or chart selectors through collisions;
  \item an object-language derivation from a successful Python or host-Lean computation alone.
\end{itemize}

When such data are supplied, the framework can often verify and propagate them.  When they are absent, the theorem should remain conditional or be reformulated.

\subsection{Metastability and rate extraction}

A metastable statement has the form
\[
 \forall p\,\forall g\,\exists N\leq B(p,g)\,
 \forall i,j\in[N,N+g(N)]\;d(a_i,a_j)\leq2^{-p}.
\]
If $B$ is primitive recursive and the displayed relation is bounded after fixing $p,g,N$, least bounded search can produce $N$.  This is an excellent iPRA statement.  It does not automatically yield a Cauchy modulus $M(p)$ independent of $g$.  Extracting such a modulus needs an additional mathematical argument.  OCiPRA treats metastability as its own observation contract rather than silently strengthening it to convergence with rate.

\subsection{Uniformity is theorem data}

A proof for each fixed dimension, degree, modulus, or number of variables need not be uniform when that parameter becomes an input.  The generated object term may grow with the externally fixed parameter.  Claims of a single primitive-recursive algorithm require a uniform syntax construction and uniform resource bound.  The report therefore distinguishes:

\begin{itemize}
  \item a family of compiled theorems indexed externally by $d$;
  \item one theorem with $d$ as an object-language variable.
\end{itemize}

This distinction is especially important for nested descent, elimination matrices, iterated differentiation, and combinatorial topology.

\subsection{Statement maximality}

For several obstruction classes, the relational or guarded theorem is plausibly the strongest generally valid constructive statement.  Examples include:

\begin{itemize}
  \item a root cluster rather than a globally labelled root;
  \item a quaternion or projective axis rather than an exact normalized axis at singular rotations;
  \item a finite cover rather than a selected chart;
  \item a sorted spectrum rather than a coherent eigenbasis;
  \item an open error interval rather than a sharp closed endpoint;
  \item a conditional complete list from a height certificate rather than a list from finiteness alone.
\end{itemize}

A future proof of a stronger statement must identify the extra information that removes the obstruction.  The admissibility classifier is designed to make that information explicit.

\subsection{Priority 1: finish the generic object compiler}

The most valuable immediate project is not another large application packet.  It is a small, reusable Lean library completing:

\begin{enumerate}
  \item bounded formula syntax and Booleanization;
  \item generic bounded least search;
  \item a fixed-DAG proof composer;
  \item a clocked trace invariant compiler;
  \item finite viability and stable-observable elimination.
\end{enumerate}

These five modules would turn many existing packet outlines into a finite sequence of local arithmetic obligations.  They would also establish a common proof-status vocabulary: semantic theorem, local object contract, globally compiled derivation, and exported certificate.

\subsection{Priority 2: consolidate existing object combinators}

The repository reports object-language results for composition, variable-length iteration, relational simulation, repeated squaring, term compilation, backward-demand equations, and stable observations.  They currently live in separate packets.  A consolidation effort should:

\begin{itemize}
  \item port the strongest versions into one root Lean namespace;
  \item deduplicate copied kernel sources;
  \item standardize arity and substitution helpers;
  \item add small cross-module examples;
  \item document exact theorem dependencies with \code{\#print axioms} output.
\end{itemize}

The OCiPRA interfaces provide a target organization for that work.

\subsection{Priority 3: an automated packet linter}

The profile classifier can be extended into a repository tool.  Given a machine-readable theorem manifest, it would check for:

\begin{itemize}
  \item a fixed or clocked computation;
  \item explicit primitive-recursive fuel;
  \item bounded certificate predicates;
  \item producer status;
  \item collision and separation fields;
  \item compactness and convergence moduli;
  \item an intended object-proof endpoint;
  \item exact status of every Lean declaration.
\end{itemize}

The linter should fail closed: an omitted field is a visible unresolved obligation, not an assumed theorem.  It cannot certify mathematical correctness, but it can prevent several recurring category errors.

\subsection{Priority 4: proof-term size and normalization}

A generic compiler may produce very large natural-deduction terms.  Research questions include:

\begin{enumerate}
  \item Which local proofs should remain native until the final transfer?
  \item Can common equality transports be shared by a DAG proof representation?
  \item How small are encoded proof trees after maximal compression?
  \item Can a verified proof-term normalizer preserve the compact checker interface?
  \item Is it better to compile a whole trace invariant or compose many local contracts directly?
\end{enumerate}

The answer may depend on whether the goal is kernel elaboration speed, standalone checker size, or certificate size.  OCiPRA separates the mathematical compiler from these optimization choices.

\subsection{Priority 5: producer theorems}

Once verifier layers are standardized, the remaining research questions become sharper:

\begin{itemize}
  \item Which classes of curves admit primitive-recursive height-certificate producers provable in iPRA?
  \item Which descent coverings admit uniform rational representative production under explicit local hypotheses?
  \item Which analytic number-theory estimates can be supplied with fully explicit constants and primitive-recursive thresholds?
  \item Which singular algebraic branches admit finite chart producers without nonconstructive root labels?
  \item Which compactness arguments can be replaced by computable finite-net moduli in the intended presentation?
\end{itemize}

These are mathematical producer problems, not bookkeeping.  A completed verifier makes their exact required output testable.

\subsection{Further mathematical questions}

The framework suggests several general investigations.

\begin{enumerate}
  \item \textbf{Minimal observation types.}  For a theorem, what is the coarsest stable observable that still implies the desired application?  Replacing representatives by invariants may drastically simplify formalization.
  \item \textbf{Demand inference.}  Can precision transformers and finite error DAGs be synthesized from local rational inequalities and then verified in iPRA?
  \item \textbf{Relational normalization.}  Is there a normal form turning any bounded nondeterministic finite computation into backward viability plus a stable observable?
  \item \textbf{Certificate completeness.}  For which mathematical domains can every valid finite observation be accompanied by a bounded certificate whose checker has a small object proof?
  \item \textbf{Negative calibration.}  Can omniscience reductions be formalized systematically to prove that proposed exact selectors exceed the intended constructive strength?
  \item \textbf{Bi-interpretation-aware compilation.}  Which proofs are shortest in native scalar syntax, which in ordinary iPRA, and can backend selection be automated without increasing the trusted checker?
  \item \textbf{Proof-producing interval arithmetic.}  A generic rational inequality compiler could discharge most analytic error bookkeeping and return native or ordinary derivations.
  \item \textbf{Finite homotopy data.}  Develop stable observables for nerves, cubical chains, and subdivision certificates without selecting global points or paths.
\end{enumerate}

\subsection{Falsifiability criteria}

A claimed OCiPRA formalization is refuted by any of the following:

\begin{itemize}
  \item a theorem endpoint that is only a host \code{Prop};
  \item an unbounded loop or search hidden in a supposedly primitive-recursive program;
  \item a certificate producer called by the verifier without a bound or proof;
  \item an exact selector at a collision with no stability or gap theorem;
  \item a ``finite'' set enumerated without a height/coding bound;
  \item an all-precision property used directly as a quantifier-free induction invariant;
  \item a compiled target proposition whose existence is confused with a proof of it;
  \item an axiom audit containing \code{sorryAx}, a custom axiom, or an unintended choice principle.
\end{itemize}

These criteria make the framework operational rather than rhetorical.

\appendix
\section{Packet-conversion workflow}

This appendix gives a concrete procedure for converting a theorem packet into an OCiPRA development.

\subsection{Step 1: normalize the theorem statement}

Identify the externally quantified numerical inputs and the requested observation.  Replace ideal higher-type objects by finite presentations.  Decide whether the output is:

\begin{enumerate}
  \item deterministic;
  \item bounded relational;
  \item a stable observable of a relation;
  \item conditional on a producer certificate.
\end{enumerate}

Record every exact selector and ask whether collisions or boundary cases are possible.  Add a gap hypothesis or weaken the output before coding begins.

\subsection{Step 2: draw the contract graph}

For every node, record:

\begin{itemize}
  \item fixed primitive-recursive program or bounded relation;
  \item input and output code conventions;
  \item precondition and postcondition;
  \item precision demand transformer;
  \item resource bound;
  \item producer data;
  \item local object-proof endpoint;
  \item intended backend: native, full iPRA, or mixed.
\end{itemize}

Topologically sort deterministic dependencies.  For adaptive or input-coded computations, specify the fixed transition system and clock.

\subsection{Step 3: make every verifier bounded}

Rewrite certificate validity using bounded quantifiers and fixed primitive-recursive terms.  Supply explicit bounds for list lengths, coefficient ranges, matrix dimensions, denominators, trace lengths, and search boxes.  A malformed code receives a default interpretation and fails the validity bit; partial decoding is not allowed.

\subsection{Step 4: isolate producers}

For every modulus, height, finite cover, representative, exceptional zero, chart, or analytic constant, choose one of:

\begin{itemize}
  \item a proved primitive-recursive producer;
  \item an explicit theorem input with a bounded verification relation;
  \item an unresolved target, visibly preventing the unconditional theorem.
\end{itemize}

Do not place producer code inside the verifier unless its totality and resource bound are already proved.

\subsection{Step 5: prove local contracts}

Generate equational proofs in \code{NPrf} where possible.  Use ordinary \code{Prf PRAAx} for bounded logical structure, witnesses, and implications.  Every host theorem used to justify a local contract must be followed by an actual object derivation with the same statement.

\subsection{Step 6: compile and audit}

Run reverse demand and resource analysis.  Compile deterministic nodes, trace loops, and relational viability.  Apply stability only when its proof is present.  Then audit the result using the gates of \cref{sec:artifact-map} and the theorem ledger.

\section{Reusable theorem schemas}

\subsection{Deterministic schema}

Data:
\[
 P(x),\quad f(x),\quad Q(y,x),\quad
 P(x)\vdash Q(f(x),x).
\]
Conclusion:
\[
 \vdash P(x)\to\exists y\,Q(y,x).
\]

\subsection{Bounded relational schema}

Data:
\[
 P(x),\quad B(x),\quad R(y,x),\quad
 P(x)\vdash\exists y\leq B(x)\,R(y,x).
\]
Conclusion: a least witness program and proof of validity/minimality, or the relation itself when no canonical witness is needed.

\subsection{Stable-observable schema}

Data: a bounded successful-transcript relation $S(x,\tau)$, a proof of existence, and
\[
 S(x,\tau)\wedge S(x,\sigma)\vdash O(\tau)=O(\sigma).
\]
Conclusion: a deterministic primitive-recursive observation independent of the selected transcript.

\subsection{Quantitative schema}

Replace equality by
\[
 d(O(\tau),O(\sigma))\leq\varepsilon(p).
\]
Conclusion: an approximate observation with the certified diameter; a regular name follows from cross-precision coherence and a shrinking error budget.

\subsection{Conditional producer schema}

Data:
\[
 \mathsf{Cert}(x,c),\qquad
 \mathsf{Cert}(x,c)\vdash\Phi(x,c).
\]
Conclusion:
\[
 \vdash (\exists c\,\mathsf{Cert}(x,c))\to\exists c\,\Phi(x,c).
\]
An unconditional theorem additionally requires a producer proof.

\section{Machine-readable contract schema}

The archive contains \path{schema/contract.schema.json}.  Its intended fields include:

\begin{itemize}
  \item theorem name and statement class;
  \item fixed/clocked computation status;
  \item primitive-recursive resource and demand descriptions;
  \item bounded verifier status;
  \item producer status;
  \item collision and separation information;
  \item compactness, convergence, height, and constant certificates;
  \item local object-derivation status;
  \item final Lean endpoint and audit status.
\end{itemize}

A manifest conforming to the schema is documentation, not a proof.  Its value is to make omitted boundaries machine-detectable.

\section{Artifact map and reproduction}\label{sec:artifact-map}

The archive is organized as follows.

\begin{longtable}{@{}p{0.31\textwidth}p{0.65\textwidth}@{}}
\toprule
Path & Contents\\
\midrule
\endfirsthead
\toprule
Path & Contents\\
\midrule
\endhead
\path{paper/ocipra_framework.tex} & Main LaTeX source.\\
\path{paper/sections/} & Section sources containing all human arguments.\\
\path{paper/ocipra_framework.pdf} & Compiled report.\\
\path{lean/OCiPRA/Contracts.lean} & Proof-carrying object interfaces and native-to-PRA transfer.\\
\path{lean/OCiPRA/TargetInterfaces.lean} & Exact packages whose proof fields remain to be instantiated.\\
\path{lean/OCiPRA/CompilerBlueprint.lean} & Backend design and target compiler signatures.\\
\path{lean/OCiPRA/Examples.lean} & Small actual object-language identity contract.\\
\path{lean/FORMALIZATION_PLAN.md} & Dependency-ordered completion plan.\\
\path{code/ocipra.py} & Exact finite reference algorithms.\\
\path{code/test_ocipra.py} & Eleven regression tests.\\
\path{code/demo.py} & Reproducible examples and JSON output.\\
\path{schema/contract.schema.json} & Machine-readable theorem-profile schema.\\
\path{metadata/THEOREM_LEDGER.md} & Claim-by-claim proof and implementation status.\\
\path{metadata/DIFFICULTY_MATRIX.csv} & Repository obstruction map.\\
\path{provenance/} & Pinned source inventory and audit notes.\\
\path{validation/} & Test, PDF build, metadata, and visual-inspection records.\\
\bottomrule
\end{longtable}

Run from the archive root:

\begin{lstlisting}
python3 verify.py
bash build.sh
\end{lstlisting}

The verification command runs the Python unit tests and exact examples.  The build script compiles the PDF when a standard \code{pdflatex} is available and refreshes the manifest.  Lean is intentionally not downloaded by the script.  To test the Lean outline, place it in the pinned repository workspace with Lean 4.32.0 and resolve the local module path as described in \path{lean/FORMALIZATION_PLAN.md}.

\section{Proof-status conventions}

The archive uses the following labels.

\begin{description}
  \item[Human proof complete.] The report gives a detailed mathematical argument with finite data and induction audit.
  \item[Executable finite check.] The Python program replays finite instances using exact arithmetic.  This is not a universal proof.
  \item[Lean interface only.] Definitions and exact target signatures exist, but no kernel build or proof is claimed.
  \item[Object derivation.] A declaration's conclusion is an actual \code{NPrf} or \code{Prf PRAAx} judgment.
  \item[Producer open.] The verifier theorem is developed, but a uniform certificate-producing theorem is not supplied.
  \item[Reformulated.] The original exact statement is replaced by a relational, approximate, quotient-level, or guarded form because the unrestricted selector is constructively invalid or unsupported.
\end{description}

The labels are intentionally asymmetric.  Compiling a Lean definition of a target proposition is not evidence that the proposition has been proved.  Passing finite tests is evidence about an implementation, not about all inputs.  A complete human proof is not advertised as kernel checked.

\section{Source baseline}

The primary source is the \repo{} repository at commit
\path{9df5cbdf39e6220be823f611987bd584d37711fd}.  The source audit focused on:

\begin{itemize}
  \item the root primitive-recursive arithmetic status pages;
  \item \code{ConciseIPRA} syntax, deduction, primitive-recursive, native, and bi-interpretation modules;
  \item the foundation packets on finite observations, finite demand, contracts, relational composition, resources, local-to-global compilation, coherent methods, and presentations-to-witnesses;
  \item selected application difficulty reports for algebra, geometry, probability, Diophantine geometry, descent, topology, and analytic number theory.
\end{itemize}

The audit is targeted, not a re-verification of all seventy-two packets.  Existing contributions are credited as inputs; the present report's addition is their proof-producing synthesis, finite demand-closure theorem, statement calculus, and conversion workflow.

\section{Conclusion}

The central lesson is simple but strict.  Mathematics fits iPRA when every purported infinity is accessed through a finite observation, every computation has primitive-recursive resources, every choice is bounded or relational, every nonuniform datum is exposed as a producer certificate, and every semantic argument ends in an actual object derivation.  Under those conditions, diverse proofs reduce to a small collection of compilers: Booleanization, bounded search, contract composition, clocked trace induction, finite demand closure, backward viability, and stable-observable elimination.

The existing repository already contains substantial pieces of this program.  OCiPRA supplies a common architecture and a concrete completion order.  Its strongest practical recommendation is to finish the generic object-language compiler before adding further large application outlines.  That compiler would turn the collection from many parallel finitary arguments into a reusable proof-producing library.


\begin{thebibliography}{99}
\bibitem{ProveIt}
V. Reshetnikov, \emph{ProveIt}, commit
\code{9df5cbdf39e6220be823f611987bd584d37711fd}, 2026.
\url{https://github.com/VladimirReshetnikov/ProveIt}

\bibitem{ConciseIPRA}
V. Reshetnikov, \emph{ConciseIPRA Lean library}, in \repo{}, files
\code{PR.lean}, \code{Deduction.lean}, \code{Native.lean}, and
\code{NativeBiInterpretation.lean}, 2026.

\bibitem{BishopBridges}
E. Bishop and D. Bridges, \emph{Constructive Analysis}, Springer, 1985.

\bibitem{BridgesRichman}
D. Bridges and F. Richman, \emph{Varieties of Constructive Mathematics}, Cambridge University Press, 1987.

\bibitem{TroelstraVanDalen}
A. S. Troelstra and D. van Dalen, \emph{Constructivism in Mathematics}, Volumes I--II, North-Holland, 1988.

\bibitem{Kohlenbach}
U. Kohlenbach, \emph{Applied Proof Theory: Proof Interpretations and their Use in Mathematics}, Springer, 2008.

\bibitem{Kleene}
S. C. Kleene, \emph{Introduction to Metamathematics}, North-Holland, 1952.
\end{thebibliography}

\end{document}
